\documentclass[11pt,a4paper]{article}
\usepackage[T2A]{fontenc}
\usepackage[utf8]{inputenc}
\usepackage[russian]{babel}
\usepackage{amsmath,amssymb,amsthm,mathtools}
\usepackage[margin=2.3cm]{geometry}
\usepackage{enumitem}
\usepackage{tikz}
\usetikzlibrary{arrows.meta,calc,patterns,decorations.pathreplacing}
\usepackage{tikz-cd}
\usepackage[pdfencoding=auto,psdextra,colorlinks=true,linkcolor=blue!50!black,citecolor=green!40!black,urlcolor=blue!60!black]{hyperref}

\setlist{itemsep=2pt,topsep=4pt}

\theoremstyle{plain}
\newtheorem{theorem}{Теорема}[section]
\newtheorem{proposition}[theorem]{Предложение}
\newtheorem{lemma}[theorem]{Лемма}
\newtheorem{corollary}[theorem]{Следствие}
\theoremstyle{definition}
\newtheorem{definition}[theorem]{Определение}
\newtheorem{example}[theorem]{Пример}
\theoremstyle{remark}
\newtheorem{remark}[theorem]{Замечание}

\renewcommand{\C}{\mathbb{C}}
\newcommand{\Z}{\mathbb{Z}}
\newcommand{\Q}{\mathbb{Q}}
\newcommand{\R}{\mathbb{R}}
\newcommand{\PP}{\mathbb{P}}
\newcommand{\A}{\mathbb{A}}
\newcommand{\cO}{\mathcal{O}}
\newcommand{\cL}{\mathcal{L}}
\newcommand{\cE}{\mathcal{E}}
\newcommand{\cF}{\mathcal{F}}
\newcommand{\cK}{\mathcal{K}}
\newcommand{\cI}{\mathcal{I}}
\newcommand{\mm}{\mathfrak{m}}
\newcommand{\pp}{\mathfrak{p}}
\DeclareMathOperator{\Pic}{Pic}
\DeclareMathOperator{\Cl}{Cl}
\DeclareMathOperator{\Div}{Div}
\DeclareMathOperator{\CaDiv}{CaDiv}
\DeclareMathOperator{\CaCl}{CaCl}
\DeclareMathOperator{\divv}{div}
\DeclareMathOperator{\ord}{ord}
\DeclareMathOperator{\Gr}{Gr}
\DeclareMathOperator{\Bl}{Bl}
\DeclareMathOperator{\Hom}{Hom}
\DeclareMathOperator{\rk}{rk}
\DeclareMathOperator{\codim}{codim}
\DeclareMathOperator{\Sing}{Sing}
\DeclareMathOperator{\Res}{Res}
\DeclareMathOperator{\GL}{GL}
\DeclareMathOperator{\PGL}{PGL}
\DeclareMathOperator{\Cone}{Cone}
\DeclareMathOperator{\conv}{conv}
\DeclareMathOperator{\NE}{\overline{NE}}
\DeclareMathOperator{\trdeg}{tr.deg}
\DeclareMathOperator{\im}{Im}
\DeclareMathOperator{\id}{id}
\DeclareMathOperator{\Ric}{Ric}
\DeclareMathOperator{\vol}{vol}
\DeclareMathOperator{\Sym}{Sym}
\newcommand{\an}{\mathrm{an}}
\newcommand{\Zar}{\mathrm{Zar}}

\title{Дивизоры, канонический класс, многообразия Фано\\ и кубические поверхности\\[4pt]
\large Конспект по классической алгебраической геометрии над $\C$}
\author{}
\date{}

\begin{document}
\maketitle

\begin{abstract}
Конспект написан на языке классической (не схемной) алгебраической геометрии: все многообразия
квазипроективны и рассматриваются над полем $\C$. Раздел~0 содержит необходимые предварительные
сведения; самые трудные классические теоремы коммутативной алгебры и анализа (теорема Гильберта
о нулях, факториальность регулярных локальных колец, теоремы Серра, GAGA, теорема Кодаиры и~т.\,п.)
формулируются со ссылками, остальное доказывается. Разделы~1--8 отвечают на поставленные вопросы:
дивизоры и линейные расслоения, касательные пространства, канонический класс, торические
многообразия, формула присоединения, обильность и многообразия Фано, раздутия, кубические
поверхности и 27 прямых.
\end{abstract}

\tableofcontents

\setcounter{section}{-1}
% ============================================================
\section{Предварительные сведения}\label{sec:0}
% ============================================================

Всюду основное поле --- $\C$. Слово \emph{многообразие} означает неприводимое квазипроективное
алгебраическое многообразие над $\C$ (если не оговорено противное). Точка всегда означает
обычную («замкнутую») точку.

\subsection{Аффинные многообразия}

Пусть $\A^n=\C^n$ и $\C[x]=\C[x_1,\dots,x_n]$. Для множества $S\subset\C[x]$ положим
$V(S)=\{p\in\A^n: f(p)=0\ \forall f\in S\}$, а для $Y\subset \A^n$ положим
$I(Y)=\{f: f|_Y=0\}$.

\begin{definition}
Множества вида $V(S)$ называются \emph{замкнутыми}; они являются замкнутыми множествами
топологии (\emph{топологии Зарисского}) на $\A^n$, так как
$V(S)\cup V(T)=V(ST)$, $\bigcap V(S_\alpha)=V(\bigcup S_\alpha)$. Замкнутое множество
\emph{неприводимо}, если оно не является объединением двух собственных замкнутых подмножеств.
\emph{Аффинное многообразие} --- неприводимое замкнутое подмножество $X\subset\A^n$;
его \emph{координатное кольцо} --- $\C[X]=\C[x]/I(X)$.
\end{definition}

\begin{theorem}[Гильберт, теорема о нулях]\label{thm:Nullstellensatz}
Для идеала $J\subset \C[x]$ имеем $I(V(J))=\sqrt J$. В частности, $V(J)=\varnothing$ тогда и
только тогда, когда $1\in J$.
\end{theorem}
Доказательство см. \cite[гл.~I, \S\,2]{Sh}, \cite[I.1.3A]{Ha}. Следствие: соответствие
$Y\mapsto I(Y)$ --- биекция между замкнутыми множествами и радикальными идеалами; неприводимым
множествам отвечают простые идеалы, точкам --- максимальные идеалы
$\mm_p=(x_1-p_1,\dots,x_n-p_n)$. Пространство $\A^n$ \emph{нётерово}: всякое замкнутое множество
есть конечное объединение неприводимых (\emph{неприводимых компонент}).

\begin{definition}
Пусть $X$ --- аффинное многообразие, $U\subset X$ открыто. Функция $f\colon U\to \C$
\emph{регулярна}, если всякая точка $p\in U$ имеет окрестность $W$, на которой $f=g/h$, где
$g,h\in \C[X]$, $h$ не обращается в нуль на $W$. Обозначение: $\cO_X(U)$.
\end{definition}

\begin{proposition}\label{prop:OX=CX}
$\cO_X(X)=\C[X]$. Более общо, для $h\in\C[X]$, $h\ne0$, и главного открытого множества
$X_h=\{h\neq 0\}$ имеем $\cO_X(X_h)=\C[X][h^{-1}]$.
\end{proposition}
\begin{proof}
Пусть $f\in\cO_X(X_h)$. Рассмотрим идеал $J=\{a\in \C[X]: af\in \C[X]\}$. По определению каждая
точка $p\in X_h$ имеет представление $f=g/b$ с $b(p)\ne 0$, и тогда $b\in J$; значит,
$V(J)\subset V(h)$. По теореме о нулях $h^N\in J$ для некоторого $N$, т.\,е. $f=a/h^N$ с
$a\in\C[X]$.
\end{proof}

\emph{Морфизм} аффинных многообразий $X\subset\A^n\to Y\subset \A^m$ --- отображение, заданное
многочленами; он задаёт гомоморфизм $\varphi^*\colon \C[Y]\to\C[X]$, и это даёт
антиэквивалентность категории аффинных многообразий и категории конечно порождённых
$\C$-алгебр без делителей нуля \cite[I.3]{Ha}.

\subsection{Проективные и квазипроективные многообразия}

$\PP^n$ --- множество прямых в $\C^{n+1}$, точки записываются в однородных координатах
$[x_0:\dots:x_n]$. Для однородных многочленов $F$ корректно определено множество нулей; замкнутые
множества $\PP^n$ --- это $V_+(S)$ для наборов однородных многочленов $S$. Покрытие
$\PP^n=\bigcup U_i$, $U_i=\{x_i\neq 0\}\cong \A^n$ (координаты $x_j/x_i$), задаёт
\emph{стандартные аффинные карты}.

\begin{definition}
\emph{Проективное многообразие} --- неприводимое замкнутое подмножество $\PP^n$;
\emph{квазипроективное многообразие} --- открытое подмножество проективного. Функция на открытом
$U\subset X$ регулярна, если её ограничения на $U\cap U_i$ регулярны в смысле аффинных многообразий.
Морфизм $\varphi\colon X\to Y$ --- непрерывное отображение, при котором прообраз регулярной функции
регулярен. Так возникает пучок колец $\cO_X$ (см.~\S\ref{ss:sheaves}).
\end{definition}

Примеры морфизмов: \emph{вложение Сегре} $\PP^n\times\PP^m\to\PP^{(n+1)(m+1)-1}$,
$([x],[y])\mapsto [x_iy_j]$; образ замкнут, что делает произведение проективных многообразий
проективным многообразием. \emph{Вложение Веронезе} $v_d\colon\PP^n\to\PP^{\binom{n+d}{d}-1}$
задаётся всеми мономами степени $d$.

\begin{theorem}[{\cite[I.5]{Sh}, \cite[I.3.4, II.4.9]{Ha}}]\label{thm:complete}
Пусть $X$ --- проективное многообразие. Тогда (а) $\cO_X(X)=\C$; (б) для любого морфизма
$\varphi\colon X\to Y$ в квазипроективное многообразие образ $\varphi(X)$ замкнут.
\end{theorem}

\subsection{Рациональные функции, размерность, бирациональность}

\emph{Поле рациональных функций} $\C(X)$ аффинного многообразия --- поле частных $\C[X]$; для
квазипроективного $X$ это $\C(X\cap U_i)$ для любой аффинной карты, пересекающей $X$. Рациональная
функция $f$ \emph{регулярна в точке} $p$, если $f=g/h$ около $p$ с $h(p)\ne0$.
\emph{Размерность} $\dim X=\trdeg_\C \C(X)$; она совпадает с максимальной длиной цепочки
неприводимых замкнутых подмножеств. \emph{Рациональное отображение} $X\dashrightarrow Y$ ---
морфизм, определённый на непустом открытом подмножестве (с отождествлением совпадающих на
открытом); $X$ и $Y$ \emph{бирационально эквивалентны}, если есть взаимно обратные рациональные
отображения; это равносильно $\C(X)\cong\C(Y)$. Многообразие \emph{рационально}, если оно
бирационально $\PP^n$.

\begin{theorem}[{Крулль, о главных идеалах; \cite[I.6, Теорема 5]{Sh}, \cite[I.1.11A]{Ha}}]\label{thm:Krull}
Если $f\in\C[X]$ не равна нулю тождественно и не обратима, то все неприводимые компоненты $V(f)$
имеют размерность $\dim X-1$.
\end{theorem}

\subsection{Локальные кольца, гладкие и нормальные точки}

\begin{definition}
\emph{Локальное кольцо} точки $p\in X$ --- кольцо $\cO_{X,p}\subset\C(X)$ рациональных функций,
регулярных в $p$; это локальное кольцо с максимальным идеалом $\mm_p=\{f: f(p)=0\}$ и
$\cO_{X,p}/\mm_p=\C$. Если $X$ аффинно, то $\cO_{X,p}=\C[X]_{\mm_p}$ (локализация).
Если $Y\subset X$ --- неприводимое замкнутое подмножество, то $\cO_{X,Y}\subset \C(X)$ ---
кольцо функций, регулярных хотя бы в одной точке $Y$ (локализация $\C[X]$ по простому идеалу
$I(Y)$ в аффинной карте).
\end{definition}

\begin{definition}
Точка $p\in X$ \emph{гладкая} (неособая), если $\dim_\C \mm_p/\mm_p^2=\dim X$. Многообразие
\emph{гладкое}, если все его точки гладкие.
\end{definition}

В~\S\ref{sec:tangent} мы увидим, что $(\mm_p/\mm_p^2)^*$ --- касательное пространство Зарисского и
что всегда $\dim \mm_p/\mm_p^2\ge \dim X$, причём множество особых точек $\Sing X$ --- собственное
замкнутое подмножество. Если $p$ гладкая, то $\cO_{X,p}$ --- \emph{регулярное локальное кольцо}.

\begin{theorem}[{Ауслендер--Буксбаум; \cite[II.3, Теорема 1]{Sh}, \cite[II.6.11.1A]{Ha}}]\label{thm:UFD}
Регулярное локальное кольцо факториально. В частности, локальное кольцо гладкой точки
многообразия --- область с однозначным разложением на множители.
\end{theorem}

\begin{definition}
$X$ \emph{нормально} в точке $p$, если $\cO_{X,p}$ целозамкнуто в $\C(X)$; $X$ нормально, если
нормально во всех точках. $X$ \emph{локально факториально}, если все $\cO_{X,p}$ факториальны.
\end{definition}

\begin{lemma}
Факториальное кольцо $A$ целозамкнуто. Поэтому гладкое многообразие нормально.
\end{lemma}
\begin{proof}
Пусть $x=a/b$ ($a,b\in A$ взаимно просты) удовлетворяет
$x^m+c_1x^{m-1}+\dots+c_m=0$, $c_i\in A$. Умножая на $b^m$, получаем
$a^m=-b(c_1a^{m-1}+\dots+c_mb^{m-1})$. Если $b$ необратим, то простой делитель $b$ делит $a^m$,
а значит, и $a$ --- противоречие. Итак, $x\in A$.
\end{proof}

\begin{theorem}[{\cite[I.6.2A, II.6.3A]{Ha}, \cite[II.5]{Sh}}]\label{thm:normal}
Пусть $X$ нормально.
\begin{enumerate}
\item Для любого неприводимого замкнутого $Y\subset X$ коразмерности~1 (\emph{простого дивизора})
кольцо $\cO_{X,Y}$ --- кольцо дискретного нормирования. Соответствующее нормирование
$\ord_Y\colon\C(X)^*\to\Z$ называется \emph{порядком} функции вдоль $Y$.
\item (\emph{Алгебраическая теорема Хартогса.}) Если $U\subset X$ открыто, $f\in\C(X)^*$ и
$\ord_Y(f)\ge0$ для всех простых дивизоров $Y$, пересекающих $U$, то $f\in\cO_X(U)$.
\end{enumerate}
\end{theorem}

Если $Y$ содержит гладкую точку $p$ многообразия $X$ и $\pi\in\mm_p$ --- локальное уравнение $Y$
около $p$ (см.~\S\ref{sec:div}), то $\ord_Y(f)$ --- просто наибольшая степень $\pi$, делящая $f$ в
факториальном кольце $\cO_{X,p}$.

\subsection{Пучки}\label{ss:sheaves}

\begin{definition}
\emph{Предпучок} абелевых групп $\cF$ на топологическом пространстве $X$ сопоставляет каждому
открытому $U$ группу $\cF(U)$ (её элементы --- \emph{сечения} над $U$), а каждому вложению
$V\subset U$ --- гомоморфизм ограничения $s\mapsto s|_V$, причём $(s|_V)|_W=s|_W$. Предпучок
--- \emph{пучок}, если для любого покрытия $U=\bigcup U_\alpha$:
(а) если $s|_{U_\alpha}=0$ для всех $\alpha$, то $s=0$; (б) если $s_\alpha\in\cF(U_\alpha)$
согласованы на пересечениях, то существует $s\in \cF(U)$ с $s|_{U_\alpha}=s_\alpha$.
\emph{Слой} $\cF_p=\varinjlim_{U\ni p}\cF(U)$ --- ростки сечений в $p$.
\end{definition}

Основной пример --- \emph{структурный пучок} $\cO_X$ регулярных функций на многообразии $X$
с топологией Зарисского; его слои --- локальные кольца $\cO_{X,p}$. Также: пучок обратимых функций
$\cO_X^*$ (мультипликативная группа), постоянный пучок $\cK_X$ с $\cK_X(U)=\C(X)$ для непустых $U$
(это пучок, так как непустые открытые множества неприводимого $X$ попарно пересекаются), пучок
$\cK_X^*$.

\begin{definition}
\emph{Пучок $\cO_X$-модулей} $\cF$: каждое $\cF(U)$ --- $\cO_X(U)$-модуль, согласованно с
ограничениями. $\cF$ \emph{локально свободен ранга $r$}, если есть открытое покрытие
$X=\bigcup U_i$ и изоморфизмы $\cF|_{U_i}\cong\cO_{U_i}^{\oplus r}$; \emph{обратимый пучок} ---
локально свободный пучок ранга~1.
\end{definition}

Для локально свободных пучков определены $\cE\oplus\cF$, $\cE\otimes\cF$,
$\cE^\vee=\mathcal{H}om(\cE,\cO_X)$, внешние степени $\Lambda^k\cE$ и
\emph{детерминант} $\det\cE=\Lambda^{\rk \cE}\cE$ (все --- склейкой соответствующих конструкций
линейной алгебры над $\cO_{U_i}$). Последовательность пучков $0\to\cE'\to\cE\to\cE''\to0$
\emph{точна}, если она точна на слоях.

\begin{lemma}\label{lem:det}
\begin{enumerate}
\item Если $0\to\cE'\to\cE\to\cE''\to0$ --- точная последовательность локально свободных пучков
рангов $r',r,r''$, то $\det \cE\cong\det\cE'\otimes\det\cE''$.
\item $\det(\cE\otimes\cF)\cong(\det\cE)^{\otimes\rk\cF}\otimes(\det\cF)^{\otimes\rk\cE}$;
$\det(\cE^\vee)\cong(\det\cE)^\vee$.
\end{enumerate}
\end{lemma}
\begin{proof}
(1) Определим отображение $\det\cE'\otimes\det\cE''\to\det\cE$ формулой
$(e_1'\wedge\dots\wedge e'_{r'})\otimes(\bar e_1\wedge\dots\wedge\bar e_{r''})\mapsto
e_1'\wedge\dots\wedge e'_{r'}\wedge e_1\wedge\dots\wedge e_{r''}$, где $e_j$ --- любые
локальные поднятия $\bar e_j$. От выбора поднятий результат не зависит: другое поднятие отличается
на сечение $\cE'$, а $\Lambda^{r'+1}\cE'=0$. Локально, выбрав базис $\cE'$ и дополнив его
поднятиями базиса $\cE''$ до базиса $\cE$ (это возможно над достаточно малой окрестностью, где
все три пучка свободны), видим, что отображение переводит образующую в образующую, т.\,е.
является изоморфизмом.

(2) В локальных тривиализациях функции перехода пучков --- матрицы $g_{ij}$ и $h_{ij}$; у
$\cE\otimes\cF$ это кронекерово произведение $g_{ij}\otimes h_{ij}$, а
$\det(g\otimes h)=(\det g)^{\rk h}(\det h)^{\rk g}$ (для диагонализуемых матриц это видно из
собственных значений, а такие матрицы плотны). Второе утверждение следует из
$\det({}^tg^{-1})=(\det g)^{-1}$.
\end{proof}

\subsection{Векторные расслоения}\label{ss:bundles}

\begin{definition}
\emph{Алгебраическое векторное расслоение ранга $r$} над $X$ --- многообразие (или объединение
многообразий) $E$ с морфизмом $\pi\colon E\to X$ и открытым покрытием $X=\bigcup U_i$ с
изоморфизмами $\psi_i\colon \pi^{-1}(U_i)\xrightarrow{\sim} U_i\times\C^r$ над $U_i$, такими что
$\psi_i\circ\psi_j^{-1}(x,v)=(x,g_{ij}(x)v)$, где $g_{ij}\colon U_i\cap U_j\to\GL_r(\C)$ ---
регулярные отображения (\emph{функции перехода}). Они удовлетворяют условию коцикла
$g_{ij}g_{jk}=g_{ik}$ на $U_i\cap U_j\cap U_k$, $g_{ii}=1$. При $r=1$ говорят о
\emph{линейном расслоении}.
\end{definition}

Обратно, любой коцикл $\{g_{ij}\}$ задаёт расслоение: склеиваем $U_i\times\C^r$ по
$(x,v)\sim(x,g_{ij}(x)v)$. Две системы функций перехода $g_{ij}$, $g'_{ij}$ (на одном
покрытии) задают изоморфные расслоения тогда и только тогда, когда
$g'_{ij}=h_ig_{ij}h_j^{-1}$ для некоторых $h_i\colon U_i\to \GL_r$.

\emph{Пучок сечений} расслоения $E$: $\cE(U)=\{s\colon U\to E \text{ морфизм}, \pi\circ s=\id\}$.
В тривиализациях сечение --- набор вектор-функций $s_i\colon U\cap U_i\to\C^r$ с
$s_i=g_{ij}s_j$. Пучок $\cE$ локально свободен ранга $r$.

\begin{proposition}\label{prop:bundles-sheaves}
Сопоставление $E\mapsto\cE$ задаёт биекцию между классами изоморфизма векторных расслоений
ранга $r$ и классами изоморфизма локально свободных пучков ранга $r$.
\end{proposition}
\begin{proof}
Обе стороны описываются одними и теми же данными. Для локально свободного $\cE$ с
тривиализациями $\phi_i\colon\cE|_{U_i}\xrightarrow{\sim}\cO_{U_i}^r$ композиция
$\phi_i\circ\phi_j^{-1}$ есть автоморфизм $\cO^r_{U_i\cap U_j}$, т.\,е. умножение на матрицу
$g_{ij}\in \GL_r(\cO(U_i\cap U_j))$; это коцикл, а замена тривиализаций меняет его на
$h_ig_{ij}h_j^{-1}$. Обратно, по коциклу пучок склеивается из $\cO_{U_i}^r$. Расслоению с
функциями перехода $g_{ij}$ соответствует пучок сечений с теми же функциями перехода (из
описания сечений выше). Две конструкции взаимно обратны.
\end{proof}

Поэтому мы часто не различаем расслоение и его пучок сечений. Для линейных расслоений группа
$\cO^*$ коммутативна, и классы изоморфизма образуют группу (по тензорному произведению) ---
\emph{группу Пикара} $\Pic X$; см. \S\ref{sec:div}.

\subsection{Когомологии Чеха}

Для пучка абелевых групп $\cF$ и покрытия $\mathfrak U=\{U_i\}$ положим
$C^p(\mathfrak U,\cF)=\prod_{i_0<\dots<i_p}\cF(U_{i_0\dots i_p})$ со стандартным дифференциалом;
$H^p(\mathfrak U,\cF)$ --- когомологии этого комплекса, $\check H^p(X,\cF)=
\varinjlim_{\mathfrak U}H^p(\mathfrak U,\cF)$. В частности, $H^0(X,\cF)=\cF(X)$, а
$\check H^1(X,\cO_X^*)$ --- это в точности коциклы $\{g_{ij}\}$ по модулю кограниц
$\{h_ih_j^{-1}\}$. Значит, по предыдущему пункту
\[
\Pic X\cong \check H^1(X,\cO_X^*).
\]

\begin{theorem}[{Серр; \cite[III.5.2]{Ha}}]\label{thm:Serre}
Пусть $X$ --- проективное многообразие, $\cF$ --- локально свободный (или, более общо, когерентный)
пучок, $\cO_X(1)$ --- ограничение $\cO_{\PP^N}(1)$ при вложении $X\subset \PP^N$. Тогда
пространства $H^i(X,\cF)$ конечномерны, и для $m\gg0$ пучок $\cF(m)=\cF\otimes\cO_X(1)^{\otimes m}$
порождён глобальными сечениями и $H^i(X,\cF(m))=0$ при $i>0$.
\end{theorem}

\subsection{Комплексно-аналитическая точка зрения}\label{ss:analytic}

Многообразие $X\subset\PP^n$ является также подмножеством комплексного многообразия $\PP^n$ с
классической топологией; обозначим его $X^{\an}$. Если $X$ гладко, то $X^{\an}$ --- комплексное
многообразие размерности $\dim X$ (это доказывается в \S\ref{sec:tangent}). Регулярные функции
голоморфны, алгебраические векторные расслоения задают голоморфные (те же функции перехода).

\begin{theorem}[{Чжоу; GAGA Серра \cite{Se}, \cite[гл.~1, \S\,3]{GH}}]\label{thm:GAGA}
\begin{enumerate}
\item Замкнутое аналитическое подмножество $\PP^n$ алгебраично.
\item Если $X$ проективно, то функтор $\cF\mapsto\cF^{\an}$ --- эквивалентность категорий
когерентных алгебраических и когерентных аналитических пучков, сохраняющая когомологии. В
частности, всякое голоморфное векторное расслоение на $X^{\an}$ алгебраично, всякая
мероморфная функция рациональна и $\Pic X\cong\Pic X^{\an}=H^1(X^{\an},\cO^*_{X^{\an}})$.
\end{enumerate}
\end{theorem}

\subsection{Линейные системы и отображения в проективное пространство}\label{ss:linsys}

Пусть $\cL$ --- обратимый пучок, $s_0,\dots,s_N\in H^0(X,\cL)$ --- сечения без общих нулей. В
тривиализации $\cL|_{U_i}\cong\cO_{U_i}$ сечения становятся функциями $s_k^{(i)}$, и
$\varphi(p)=[s_0^{(i)}(p):\dots:s_N^{(i)}(p)]$ не зависит от $i$ (функции на пересечении
умножаются на общий обратимый множитель $g_{ij}$). Получаем морфизм
$\varphi\colon X\to\PP^N$ с $\varphi^*\cO_{\PP^N}(1)\cong\cL$, при котором $s_k=\varphi^*x_k$.
Если сечения имеют общие нули (\emph{базисные точки}), $\varphi$ --- лишь рациональное
отображение.

\subsection{Теория пересечений на поверхностях и кривых}\label{ss:intersection}

\begin{theorem}[{\cite[V.1.1]{Ha}, \cite[I]{Be}}]\label{thm:intersection}
Пусть $S$ --- гладкая проективная поверхность. Существует единственная симметричная билинейная
форма $\Pic S\times\Pic S\to\Z$, $(D,D')\mapsto D\cdot D'$, такая, что для гладких кривых $C,C'$,
пересекающихся трансверсально, $C\cdot C'=\#(C\cap C')$. При этом:
\begin{enumerate}
\item если $C\neq C'$ --- неприводимые кривые, то $C\cdot C'=\sum_{p\in C\cap C'}(C\cdot C')_p\ge0$,
где $(C\cdot C')_p=\dim_\C\cO_{S,p}/(f,g)$ --- локальная кратность пересечения
($f,g$ --- локальные уравнения); $(C\cdot C')_p\ge \mathrm{mult}_p C\cdot\mathrm{mult}_p C'$;
\item если $C$ --- неприводимая кривая, $\cL$ --- обратимый пучок, то
$\cL\cdot C=\deg(\cL|_{C})$ (степень на нормализации $C$);
\item (Безу) на $\PP^2$: $\Pic\PP^2=\Z H$, $H^2=1$, и кривые степеней $d,e$ без общих компонент пересекаются по
$de$ точкам с учётом кратностей.
\end{enumerate}
\end{theorem}

На гладкой проективной кривой $C$ рода $g$ степень дивизора $\sum n_ip_i$ есть $\sum n_i$; степень
главного дивизора равна $0$; канонический класс имеет степень $2g-2$; для кривой $C\subset\PP^N$
степень $\cO_C(1)$ равна числу точек пересечения с общей гиперплоскостью
\cite[IV.1]{Ha}. \emph{Теорема Римана--Роха}: $h^0(D)-h^0(K_C-D)=\deg D+1-g$.

\begin{theorem}[{Римана--Роха для поверхностей, \cite[V.1.6]{Ha}}]\label{thm:RR}
Для дивизора $D$ на гладкой проективной поверхности $S$:
$\chi(\cO_S(D))=\chi(\cO_S)+\tfrac12 D\cdot(D-K_S)$.
\end{theorem}

% ============================================================
\section{Дивизоры, линейные расслоения и обратимые пучки}\label{sec:div}
% ============================================================

В этом разделе $X$ --- многообразие; для дивизоров Вейля предполагаем $X$ нормальным
(например, гладким). Главный итог:
\[
\boxed{\ \text{лин. расслоения}/{\cong}\ \overset{(1)}{=}\ \text{обратимые пучки}/{\cong}\ =\ \Pic X
\ \overset{(2)}{\cong}\ \CaCl X\ \overset{(3)}{\hookrightarrow}\ \Cl X\ }
\]
где (1) и (2) --- изоморфизмы для любого многообразия, (3) --- вложение для нормального $X$,
являющееся изоморфизмом тогда и только тогда, когда $X$ локально факториально, в частности для
гладкого $X$.

\subsection{Дивизоры Вейля}

\begin{definition}
\emph{Простой дивизор} на $X$ --- неприводимое замкнутое подмногообразие коразмерности~1.
\emph{Дивизор Вейля} --- элемент свободной абелевой группы $\Div X$, порождённой простыми
дивизорами: $D=\sum n_iY_i$ (сумма конечна). $D$ \emph{эффективен} ($D\ge0$), если все $n_i\ge0$.
\end{definition}

Пусть $X$ нормально, $f\in\C(X)^*$. Тогда $\ord_Y(f)\ne0$ лишь для конечного числа $Y$:
действительно, $f$ регулярна и обратима на непустом открытом $U$, а $X\setminus U$ содержит лишь
конечное число простых дивизоров (своих неприводимых компонент коразмерности~1). Положим
\[
\divv(f)=\sum_Y\ord_Y(f)\,Y.
\]
Так как $\ord_Y(fg)=\ord_Y f+\ord_Y g$, получаем гомоморфизм $\divv\colon\C(X)^*\to\Div X$.
Дивизоры вида $\divv f$ называются \emph{главными}; $D\sim D'$ (\emph{линейная эквивалентность}),
если $D-D'$ главный.

\begin{definition}
\emph{Группа классов дивизоров} $\Cl X=\Div X/\{\text{главные}\}$.
\end{definition}

\begin{proposition}\label{prop:excision}
Если $Z\subset X$ замкнуто, $U=X\setminus Z$, то есть точная последовательность
$\bigoplus_{Z_i}\Z\, Z_i\to\Cl X\to\Cl U\to0$, где $Z_i$ --- компоненты $Z$ коразмерности~1.
В частности, $\Cl X\cong\Cl U$, если $\codim Z\ge2$.
\end{proposition}
\begin{proof}
Ограничение $Y\mapsto Y\cap U$ сюръективно на простых дивизорах (замыкание простого дивизора $U$
--- простой дивизор $X$) и переводит главные дивизоры в главные. Если $D|_U=\divv(f)|_U$, то
$D-\divv f$ сосредоточен на $Z$, т.\,е. является комбинацией $Z_i$.
\end{proof}

\begin{example}\label{ex:ClPn}
$\Cl\A^n=0$, $\Cl\PP^n\cong\Z$, причём изоморфизм задаётся степенью, а образующая --- класс
гиперплоскости $H$.

Действительно, простой дивизор $Y\subset\A^n$ задаётся простым идеалом $\pp=I(Y)$ высоты~1 в
факториальном кольце $\C[x]$. Возьмём $0\ne f\in\pp$; некоторый неприводимый множитель $g$ многочлена
$f$ лежит в $\pp$, идеал $(g)$ прост, и $0\subsetneq (g)\subset\pp$ влечёт $(g)=\pp$. Тогда
$Y=\divv(g)$ (порядок равен $1$, так как $g$ порождает максимальный идеал в $\cO_{\A^n,Y}$).
Значит, $\Cl\A^n=0$. Аналогично простой дивизор в $\PP^n$ --- гиперповерхность $V_+(G)$,
$G$ --- неприводимый однородный многочлен степени $d$; положим $\deg Y=d$. Для рациональной
функции $f=G_1/G_2$ ($G_i$ однородны одной степени) $\divv f=\sum(\text{неприв. множители }G_1)
-\sum(\text{неприв. множители }G_2)$ имеет степень $0$. С другой стороны,
$Y-dH=\divv(G/x_0^d)$, где $H=V_+(x_0)$. Поэтому $\deg\colon\Cl\PP^n\to\Z$ --- изоморфизм.
\end{example}

\subsection{Дивизоры Картье}

\begin{definition}
\emph{Дивизор Картье} на $X$ --- класс эквивалентности наборов $\{(U_i,f_i)\}$, где
$X=\bigcup U_i$ --- открытое покрытие, $f_i\in\C(X)^*$ и $f_i/f_j\in\cO^*_X(U_i\cap U_j)$ (регулярна
и обратима на пересечении). Два набора эквивалентны, если их объединение снова удовлетворяет этому
условию. $f_i$ называется \emph{локальным уравнением} дивизора на $U_i$. Сложение:
$\{(U_i,f_i)\}+\{(V_j,g_j)\}=\{(U_i\cap V_j,f_ig_j)\}$. Группа $\CaDiv X=H^0(X,\cK^*_X/\cO^*_X)$.
Дивизор $\{(X,f)\}$ называется \emph{главным}, $\CaCl X$ --- фактор по главным.
Дивизор \emph{эффективен}, если все $f_i\in\cO_X(U_i)$.
\end{definition}

Если $X$ нормально, то дивизору Картье $D=\{(U_i,f_i)\}$ сопоставим дивизор Вейля
\[
[D]=\sum_Y \ord_Y(f_i)\,Y,\qquad\text{где $i$ таково, что $Y\cap U_i\neq\varnothing$}.
\]
Корректность: если $Y$ пересекает $U_i$ и $U_j$, то $Y\cap U_i\cap U_j\neq\varnothing$
($Y$ неприводим), а $f_i/f_j$ обратима на $U_i\cap U_j$, поэтому $\ord_Y f_i=\ord_Y f_j$.
Получаем гомоморфизм $\CaDiv X\to\Div X$, переводящий главные в главные.

\subsection{Линейные расслоения и обратимые пучки}

По \S\ref{ss:bundles} линейное расслоение задаётся коциклом $g_{ij}\in\cO^*(U_i\cap U_j)$;
его пучок сечений --- обратимый пучок с теми же функциями перехода
(предложение~\ref{prop:bundles-sheaves}). Тензорное произведение линейных расслоений
отвечает произведению коциклов, двойственное --- обратному коциклу. Поэтому:

\begin{theorem}\label{thm:Pic=H1}
Классы изоморфизма линейных расслоений, классы изоморфизма обратимых пучков и
$\check H^1(X,\cO_X^*)$ --- одна и та же группа $\Pic X$.
\end{theorem}

Мы уже доказали это (предложение~\ref{prop:bundles-sheaves}); подчеркнём явный словарь.
Сечение линейного расслоения $L$ с функциями перехода $g_{ij}$ --- это набор функций
$s_i\in \cO(U_i)$ с $s_i=g_{ij}s_j$; \emph{рациональное (мероморфное) сечение} --- набор
$s_i\in\C(X)$ с тем же условием. Обратимый пучок $\cL=$ пучок сечений $L$.

\subsection{От дивизоров Картье к обратимым пучкам}

\begin{definition}
Для дивизора Картье $D=\{(U_i,f_i)\}$ определим подпучок $\cO_X(D)\subset\cK_X$:
\[
\cO_X(D)|_{U_i}=f_i^{-1}\cO_{U_i}.
\]
Это корректно: на $U_i\cap U_j$ имеем $f_i^{-1}\cO=f_j^{-1}\cO$, так как $f_j/f_i$ обратима.
\end{definition}

Для нормального $X$: $h\in\cO_X(D)(U)$ тогда и только тогда, когда $\divv(h)+[D]\ge0$ на $U$
(по теореме Хартогса~\ref{thm:normal}, примененной к $hf_i$). В частности,
$H^0(X,\cO_X(D))=\{h\in\C(X)^*:\divv h+D\ge0\}\cup\{0\}$, и \emph{полная линейная система}
$|D|=\{\divv h+D\}$ --- проективизация этого пространства.

\begin{lemma}\label{lem:OD}
\begin{enumerate}
\item $\cO_X(D)$ --- обратимый пучок; умножение на $f_i$ даёт тривиализации
$\cO_X(D)|_{U_i}\xrightarrow{\sim}\cO_{U_i}$, функции перехода --- $g_{ij}=f_i/f_j$.
\item $\cO_X(D+D')\cong\cO_X(D)\otimes\cO_X(D')$, $\cO_X(-D)\cong\cO_X(D)^\vee$.
\item $\cO_X(D)\cong\cO_X$ тогда и только тогда, когда $D$ главный.
\end{enumerate}
\end{lemma}
\begin{proof}
(1) Сечение $s$ над $V\subset U_i$ переходит в $s_i=f_is\in\cO(V)$, и $s_i=(f_i/f_j)s_j$.

(2) Умножение $\cK\otimes\cK\to\cK$ индуцирует
$f_i^{-1}\cO\otimes f_i'^{-1}\cO\to (f_if'_i)^{-1}\cO$, это изоморфизм локально (образующая
переходит в образующую). Для двойственного: спаривание $\cO(D)\otimes\cO(-D)\to\cO$ ---
умножение; оно совершенно, так как локально это $f_i^{-1}\cO\otimes f_i\cO\to\cO$.

(3) Если $D=\{(X,f)\}$, то $\cO(D)=f^{-1}\cO_X\cong\cO_X$. Обратно, пусть
$\alpha\colon\cO_X\xrightarrow\sim\cO_X(D)$, и $g=\alpha(1)\in H^0(X,\cO(D))\subset\C(X)$.
Так как $g$ порождает $f_i^{-1}\cO_{U_i}$ на каждом $U_i$, имеем $gf_i\in\cO^*(U_i)$. Значит, наборы
$\{(U_i,f_i)\}$ и $\{(X,g^{-1})\}$ эквивалентны: $D$ --- главный дивизор $g^{-1}$.
\end{proof}

\begin{theorem}\label{thm:CaCl=Pic}
Для любого (неприводимого) многообразия $X$ отображение $D\mapsto\cO_X(D)$ индуцирует изоморфизм
групп $\CaCl X\xrightarrow{\sim}\Pic X$.
\end{theorem}
\begin{proof}
По лемме~\ref{lem:OD} это гомоморфизм с ядром, состоящим из главных дивизоров. Докажем
сюръективность. Пусть $\cL$ --- обратимый пучок, $\phi_i\colon\cL|_{U_i}\xrightarrow\sim\cO_{U_i}$
--- тривиализации над покрытием непустыми открытыми множествами, $g_{ij}=\phi_i\phi_j^{-1}\in
\cO^*(U_i\cap U_j)$. Так как $X$ неприводимо, все $U_i\cap U_j$ непусты и плотны, и регулярные
функции на них --- рациональные функции на $X$. Зафиксируем индекс $0$ и положим
\[
f_i:=g_{i0}\in\cO^*(U_i\cap U_0)\subset\C(X)^*.
\]
Из условия коцикла $g_{i0}=g_{ij}g_{j0}$ на $U_i\cap U_j\cap U_0$ получаем равенство рациональных
функций $f_i/f_j=g_{ij}$; а $g_{ij}$ регулярна и обратима на $U_i\cap U_j$. Значит,
$D=\{(U_i,f_i)\}$ --- дивизор Картье с функциями перехода $f_i/f_j=g_{ij}$, т.\,е.
$\cO_X(D)\cong\cL$.
\end{proof}

Геометрический смысл: $\cL\cong \cO(D)$, где $D$ --- \emph{дивизор рационального сечения} $s$
расслоения $L$ (в доказательстве $s$ --- сечение, равное $1$ в тривиализации над $U_0$):
$D=\divv(s)=\{(U_i, s_i)\}$. Если $s$ регулярно, то $D\ge 0$ --- это «множество нулей сечения с
кратностями». Итак,
\[
\text{дивизор } D \;\longleftrightarrow\; (\cO(D),\ s_D=1),\qquad
(\cL,s)\;\longleftrightarrow\;\divv(s),
\]
и замена рационального сечения $s\mapsto hs$, $h\in\C(X)^*$, меняет дивизор на главный.

\subsection{Дивизоры Вейля и Картье}

\begin{theorem}\label{thm:Weil=Cartier}
Пусть $X$ нормально.
\begin{enumerate}
\item Гомоморфизм $\CaDiv X\to\Div X$, $D\mapsto[D]$, инъективен; поэтому
$\Pic X\cong\CaCl X\hookrightarrow\Cl X$.
\item Если $X$ локально факториально (в частности, если $X$ гладкое), то
$\CaDiv X\cong\Div X$ и $\Pic X\cong\CaCl X\cong\Cl X$.
\item Обратно, если каждый дивизор Вейля является дивизором Картье, то $X$ локально факториально.
\end{enumerate}
\end{theorem}
\begin{proof}
(1) Пусть $[D]=0$, $D=\{(U_i,f_i)\}$. Тогда $\ord_Y f_i=0$ для всех простых $Y$,
пересекающих $U_i$. По теореме Хартогса~\ref{thm:normal}(2), применённой к $f_i$ и $f_i^{-1}$,
обе функции регулярны на $U_i$, т.\,е. $f_i\in\cO^*(U_i)$ и $D=0$. Главные дивизоры переходят в
главные, а если $[D]=\divv(f)$, то $[D-\divv f]=0$, откуда $D=\divv f$ как дивизоры Картье;
значит, и $\CaCl X\to\Cl X$ инъективно.

(2) Достаточно доказать, что каждый простой дивизор $Y$ локально задаётся одним уравнением:
у каждой точки $p$ есть окрестность $U$ и $\pi\in\C(X)^*$ с $\divv(\pi)|_U=Y\cap U$. Тогда
$\{(U,\pi)\}$ склеиваются в дивизор Картье с $[D]=Y$ (на пересечениях отношения локальных
уравнений имеют нулевой дивизор и потому обратимы по пункту (1)), а гомоморфизм, заданный на
образующих, обратен к $D\mapsto[D]$.

Если $p\notin Y$, берём $U=X\setminus Y$, $\pi=1$. Пусть $p\in Y$. Можно считать $X$ аффинным.
Идеал $\pp=I(Y)\cO_{X,p}$ --- простой идеал высоты~1 в факториальном кольце $\cO_{X,p}$.
Как в примере~\ref{ex:ClPn}, он главный: $\pp=(\pi)$, где $\pi$ --- подходящий неприводимый множитель (лежащий в $\pp$)
некоторого ненулевого элемента $\pp$. Пусть $I(Y)=(h_1,\dots,h_k)$ в $\C[X]$. Тогда в $\cO_{X,p}$
$h_j=\pi a_j$ и $\pi=\sum b_jh_j$, где $a_j,b_j,\pi$ --- рациональные функции, регулярные в $p$;
все они регулярны на некоторой окрестности $U\ni p$. На $U$ идеал функций, обращающихся в нуль на
$Y$, порождён $\pi$. Значит, $\pi$ не имеет нулей на $U$ вне $Y$ (полюсов нет), а
$\ord_Y\pi=1$, так как $\pi$ порождает максимальный идеал кольца нормирования
$\cO_{X,Y}=(\cO_{X,p})_\pp$. Итак, $\divv(\pi)|_U=Y\cap U$.

Для гладкого $X$ локальная факториальность --- теорема~\ref{thm:UFD}.

(3) Пусть $p\in X$. Для нётеровой области $A=\cO_{X,p}$ факториальность равносильна
тому, что каждый простой идеал высоты~1 главный \cite[II.6.2A]{Ha}. Простой идеал высоты~1
кольца $\cO_{X,p}$ --- это $I(Y)\cO_{X,p}$ для простого дивизора $Y\ni p$. Если $Y$ --- дивизор
Картье с локальным уравнением $\pi$ около $p$, то по теореме Хартогса
$I(Y)\cO_{X,p}=\{f:\ord_Y f\ge1,\ \ord_{Y'}f\ge0\ \forall Y'\ni p\}=\pi\cO_{X,p}$.
\end{proof}

\begin{remark}[Почему условия существенны]
Изоморфизм $\CaCl\cong\Pic$ использует лишь неприводимость $X$ (в схемной геометрии нужна
целостность). Совпадение с $\Cl$ использует нормальность (без неё $\ord_Y$ даже не определён) и
факториальность.
\end{remark}

\begin{example}[Квадратичный конус: дивизор Вейля, не являющийся дивизором Картье]\label{ex:cone}
Пусть $X=\{xy=z^2\}\subset\A^3$ (нормальная поверхность с единственной особой точкой $0$),
$\ell=\{x=z=0\}$ --- прямая на $X$.
\begin{itemize}
\item $\divv(x)=2\ell$. Действительно, на $X$ из $x=0$ следует $z=0$, так что
$V(x)=\ell$. В точках $\ell$ с $y\ne0$ функция $z$ --- локальное уравнение $\ell$ (поверхность
гладкая, $x=z^2/y$, и $(y,z)$ --- локальные координаты), а $x=z^2/y$ имеет порядок~2.
\item $\ell$ не является дивизором Картье в точке $0$. Иначе по доказательству
теоремы~\ref{thm:Weil=Cartier}(3) $I(\ell)\cO_{X,0}=(x,z)$ был бы главным идеалом $(\pi)$.
Но $\mm_0/\mm_0^2$ трёхмерно с базисом $x,y,z$ (уравнение $xy-z^2$ лежит в $\mm^2$), образ
идеала $(\pi)$ в $\mm_0/\mm_0^2$ не более чем одномерен, а образ $(x,z)$ двумерен.
\item $\Cl X\cong\Z/2$ (порождается $\ell$), $\Pic X=0$. Действительно, $X\setminus\ell=\{x\ne0\}$
изоморфно $\{(x,z):x\ne0\}$ ($y=z^2/x$), координатное кольцо $\C[x^{\pm1},z]$ факториально,
и $\Cl(X\setminus\ell)=0$. По предложению~\ref{prop:excision} $\Cl X$ порождена $\ell$,
$2\ell\sim0$, а $\ell\not\sim0$, так как главный дивизор является дивизором Картье. Наконец,
$\Pic X\hookrightarrow\Cl X$ и $\ell$ --- не Картье, так что $\Pic X=0$.
\end{itemize}
Этот же пример будет разобран торическими методами в \S\ref{sec:toric}.
\end{example}

\subsection{Проективный случай}

\begin{theorem}\label{thm:proj-div}
Пусть $X\subset\PP^N$ --- проективное многообразие, $\cO_X(1)$ --- ограничение гиперплоского
расслоения.
\begin{enumerate}
\item Каждый обратимый пучок $\cL$ изоморфен $\cO_X(A-B)$, где $A,B$ --- эффективные дивизоры
Картье, являющиеся гиперплоскими сечениями $X$ в некоторых проективных вложениях
($\cO(A)\cong\cL(m)$, $\cO(B)\cong\cO_X(m)$).
\item Если $X$ гладкое, то
\[
\{\text{лин. расслоения}\}/{\cong}\;=\;\{\text{обратимые пучки}\}/{\cong}\;=\;\Pic X\;\cong\;
\CaCl X\;\cong\;\Cl X\;\cong\;\Pic X^{\an}=H^1(X^{\an},\cO_{X^\an}^*),
\]
причём в аналитической картине дивизоры --- это формальные суммы неприводимых аналитических
гиперповерхностей (по теореме Чжоу они алгебраичны), а мероморфные функции рациональны.
\end{enumerate}
\end{theorem}
\begin{proof}
(1) По теореме Серра~\ref{thm:Serre} $\cL(m)$ и $\cO_X(m)$ порождены глобальными сечениями при
$m\gg0$, а $\cL(m)=\cL(m_0)\otimes\cO_X(m-m_0)$ при этом очень обилен (произведение порождённого сечениями и
очень обильного очень обильно, \cite[II.7.10, II.7.5]{Ha}); значит, общее сечение $\cL(m)$ (соотв.
$\cO_X(m)$) --- гиперплоское сечение в соответствующем вложении. Ненулевое сечение $s$ обратимого пучка задаёт эффективный дивизор $\divv(s)$ с
$\cO(\divv s)\cong$ этому пучку. Тогда $\cL\cong\cL(m)\otimes\cO_X(m)^\vee\cong\cO(A-B)$.
(2) --- теоремы~\ref{thm:CaCl=Pic}, \ref{thm:Weil=Cartier}, \ref{thm:GAGA}.
\end{proof}

\begin{example}
$\Pic\PP^n=\Z\cdot[\cO(1)]$, $\cO(D)\cong\cO(d)$ для $\deg D=d$. Здесь $\cO(d)$ задаётся
функциями перехода $g_{ij}=(x_j/x_i)^{d}$ на $U_i\cap U_j$ (это функции перехода дивизора
$d\cdot V_+(x_0)$ с локальными уравнениями $f_i=(x_0/x_i)^d$: $f_i/f_j=(x_j/x_i)^d$);
сечения $\cO(d)$ --- однородные многочлены степени $d$. Расслоение $\cO(-1)$ ---
\emph{тавтологическое}: $\{([x],v)\in\PP^n\times\C^{n+1}: v\in\C x\}$.
\end{example}

\begin{remark}[Первый класс Чженя]
Экспоненциальная последовательность $0\to\Z\to\cO_{X^\an}\xrightarrow{\exp(2\pi i\,\cdot)}
\cO^*_{X^\an}\to0$ даёт гомоморфизм $c_1\colon\Pic X\to H^2(X,\Z)$. Для гладкой проективной
поверхности $D\cdot D'=\int_X c_1(D)\cup c_1(D')$; на кривой $\deg L=\int_C c_1(L)$.
\end{remark}

% ============================================================
\section{Касательное пространство}\label{sec:tangent}
% ============================================================

\subsection{Касательное пространство Зарисского}

\begin{definition}
\emph{Касательное пространство Зарисского} к многообразию $X$ в точке $p$:
\[
T_pX=T^{\Zar}_pX=(\mm_p/\mm_p^2)^*,\qquad \mm_p\subset\cO_{X,p}.
\]
Пространство $\mm_p/\mm_p^2$ называется \emph{кокасательным}; для $f\in\cO_{X,p}$ класс
$d_pf=[f-f(p)]\in\mm_p/\mm_p^2$ --- \emph{дифференциал} $f$ в точке $p$.
\end{definition}

\begin{proposition}\label{prop:derivations}
$T_pX$ канонически изоморфно пространству \emph{дифференцирований в точке $p$}, т.\,е.
линейных отображений $\delta\colon\cO_{X,p}\to\C$ с $\delta(fg)=f(p)\delta g+g(p)\delta f$.
\end{proposition}
\begin{proof}
Дифференцирование зануляет константы ($\delta(1)=2\delta(1)$) и $\mm_p^2$ (по правилу Лейбница),
поэтому задаёт функционал на $\mm_p/\mm_p^2$ через разложение $\cO_{X,p}=\C\oplus\mm_p$.
Обратно, функционал $\lambda$ на $\mm_p/\mm_p^2$ задаёт $\delta(f)=\lambda(d_pf)$; правило
Лейбница следует из $fg-f(p)g(p)\equiv f(p)(g-g(p))+g(p)(f-f(p))\pmod{\mm_p^2}$.
\end{proof}

\begin{theorem}[Вложенное описание]\label{thm:jacobian}
Пусть $X\subset\A^n$ --- аффинное многообразие, $I(X)=(f_1,\dots,f_m)$, $p\in X$. Тогда
\[
T_pX\cong\ker J(p)=\Big\{v\in\C^n:\ \sum_{i}\frac{\partial f_j}{\partial x_i}(p)\,v_i=0,\
j=1,\dots,m\Big\},\qquad J=\Big(\frac{\partial f_j}{\partial x_i}\Big).
\]
\end{theorem}
\begin{proof}
Пусть $M=(x_1-p_1,\dots,x_n-p_n)\subset\C[x]$. Отображение $M\to(\C^n)^*$,
$g\mapsto d_pg=\sum\frac{\partial g}{\partial x_i}(p)(x_i-p_i)$ (линейная часть ряда Тейлора)
сюръективно с ядром $M^2$, так что $M/M^2\cong(\C^n)^*$. Так как локализация не меняет $\mm/\mm^2$,
\[
\mm_p/\mm_p^2=\bar M/\bar M^2=M/(M^2+I(X))\cong(\C^n)^*/\langle d_pf_1,\dots,d_pf_m\rangle
\]
(здесь $\bar M$ --- образ $M$ в $\C[X]$; образ $I(X)$ в $M/M^2$ натянут на $d_pf_j$, так как
$d_p(af_j)=a(p)d_pf_j$ при $f_j(p)=0$). Двойственное пространство --- аннулятор
$\langle d_pf_j\rangle$ в $\C^n$, т.\,е. $\ker J(p)$.
\end{proof}

\begin{remark}
Важно брать \emph{все} образующие идеала $I(X)$, а не произвольные уравнения, задающие $X$ как
множество: для $X=\{x=0\}\subset\A^2$, заданного уравнением $x^2=0$, ядро якобиана было бы всем
$\C^2$.
\end{remark}

Для проективного $X\subset\PP^n$ можно работать в аффинной карте; для гиперповерхности
$X=V_+(F)$ (с $F$ без кратных множителей) удобно \emph{проективное касательное пространство}
$\overline{T}_pX=\{\sum_i\frac{\partial F}{\partial x_i}(p)x_i=0\}\subset\PP^n$; по формуле Эйлера
$\sum x_i\partial_iF=(\deg F)F$ оно содержит $p$ и в карте $U_0$ является замыканием аффинного
касательного пространства, сдвинутого в $p$.

\begin{theorem}[{\cite[II.1.4]{Sh}, \cite[I.5.3]{Ha}}]\label{thm:dimT}
Для всех $p\in X$ выполнено $\dim T_pX\ge\dim X$, и множество гладких точек
$\{p:\dim T_pX=\dim X\}$ открыто и непусто.
\end{theorem}
\begin{proof}[Набросок]
Функция $p\mapsto\dim\ker J(p)=n-\rk J(p)$ полунепрерывна сверху (ранг падает на замкнутом
множестве, где обращаются в нуль миноры); значит, минимум $r$ достигается на открытом множестве,
и $\dim T_pX\ge r$ всюду. Остаётся показать, что $r=\dim X$; для этого $X$ бирационально
гиперповерхности $\{g=0\}\subset\A^{d+1}$, $d=\dim X$ (теорема о примитивном элементе), а
у гиперповерхности (без кратных множителей) $\partial g/\partial x_i$ не делятся все на $g$, так
что $\rk J=1$ на непустом открытом множестве, т.\,е. $\dim T_p=d$.
\end{proof}

\begin{corollary}[Якобиев критерий]
Точка $p$ гладкая $\iff\rk J(p)=n-\dim X$.
\end{corollary}

\begin{example}
\begin{enumerate}
\item Кривая $y^2=x^3$ (каспидальная): $J=(-3x^2,\ 2y)$, $J(0)=0$, $T_0=\C^2$ --- особая точка.
\item Узловая кубика $y^2=x^2(x+1)$: $T_0=\C^2$.
\item Конус $xy=z^2$: $J(0)=0$, $T_0=\C^3$, хотя $\dim X=2$.
\item Гладкая проективная гиперповерхность $V_+(F)$: все $\partial F/\partial x_i$ не обращаются
одновременно в нуль на $X$.
\end{enumerate}
\end{example}

\emph{Дифференциал морфизма} $\varphi\colon X\to Y$, $q=\varphi(p)$:
$\varphi^*\colon\mm_q\to\mm_p$ индуцирует $\mm_q/\mm_q^2\to\mm_p/\mm_p^2$, двойственное к
которому --- $d_p\varphi\colon T_pX\to T_qY$; во вложенном описании это умножение на матрицу Якоби.

\subsection{Касательное пространство в смысле дифференциальной геометрии}

\begin{theorem}\label{thm:smooth-manifold}
Пусть $X\subset\A^n$ --- аффинное многообразие размерности $d$, $p\in X$ --- гладкая точка.
Тогда существует окрестность $W\ni p$ в $\C^n$ (в классической топологии) такая, что $X\cap W$
--- комплексное подмногообразие $W$ размерности $d$, и его голоморфное касательное
пространство в $p$ совпадает с $\ker J(p)=T^{\Zar}_pX$. В частности, если $X$ --- гладкое
квазипроективное многообразие, то $X^{\an}$ --- комплексное многообразие размерности $d$.
\end{theorem}
\begin{proof}
Поскольку $\rk J(p)=n-d$, можно выбрать $f_1,\dots,f_{n-d}\in I(X)$ с линейно независимыми
дифференциалами в $p$. Пусть $Y=V(f_1,\dots,f_{n-d})\supset X$. Ядро якобиана $Y$ в $p$
содержится в $\ker(d_pf_1,\dots,d_pf_{n-d})$ и потому имеет размерность $\le d$. Применяя
теорему~\ref{thm:Krull} $n-d$ раз, видим, что все компоненты $Y$ имеют размерность $\ge d$; для каждой
компоненты $Y'\ni p$ по теореме~\ref{thm:dimT} $\dim Y'\le\dim T_pY'\le\dim T_pY\le d$. Значит,
все компоненты через $p$ имеют размерность $d$, и $\dim T_pY=d$, т.\,е. $Y$ гладко в $p$. Локальное кольцо $\cO_{Y,p}$ регулярно,
значит, является областью (теорема~\ref{thm:UFD}), поэтому через $p$ проходит
\emph{единственная} неприводимая компонента $Y_0$ многообразия $Y$, и $\dim Y_0=d$. Так как
$X\subset Y$ неприводимо, содержит $p$ и имеет размерность $d$, получаем $X=Y_0$. Остальные
компоненты $Y$ замкнуты и не содержат $p$, поэтому в некоторой окрестности $W$ точки $p$
имеем $X\cap W=Y\cap W=\{f_1=\dots=f_{n-d}=0\}\cap W$. По голоморфной теореме о неявной
функции это множество --- комплексное подмногообразие размерности $d$: после перенумерации
координат $X\cap W$ --- график голоморфного отображения
$\psi\colon(x_1,\dots,x_d)\mapsto(x_{d+1},\dots,x_n)$.

Голоморфное касательное пространство подмногообразия --- это образ $d\Psi$, где
$\Psi(x')=(x',\psi(x'))$, т.\,е. пространство скоростей $\gamma'(0)$ голоморфных кривых
$\gamma\subset X$, $\gamma(0)=p$. Для $f\in I(X)$ имеем $f\circ\gamma\equiv0$, откуда
$J(p)\gamma'(0)=0$; итак, $T_pX^{\an}\subset\ker J(p)$, а размерности обоих равны $d$.

Наконец, карты $X\cap W\to\C^d$ (проекции на подходящие $d$ координат) голоморфно согласованы,
так что $X^{\an}$ --- комплексное многообразие; для квазипроективного $X$ применяем это в
аффинных картах.
\end{proof}

\begin{remark}[Вещественная точка зрения]
Как гладкое вещественное многообразие $X^{\an}$ имеет размерность $2d$, и его касательное
пространство $T^{\R}_pX^{\an}\subset\R^{2n}=\C^n$ --- это график $d\psi$, т.\,е. $\C$-линейное
подпространство: оно инвариантно относительно комплексной структуры $J_0$ (умножения на $i$).
Как $\C$-векторное пространство $(T_p^\R X^{\an},J_0)$ совпадает с $\ker J(p)$; а в
комплексифицированном касательном пространстве $T^\R_p\otimes\C=T^{1,0}\oplus T^{0,1}$
голоморфное касательное пространство --- это $T^{1,0}$, изоморфное $(T^\R_p,J_0)$ по
$v\mapsto\frac12(v-iJ_0v)$. Итак, все три пространства --- \emph{одно и то же}.
\end{remark}

\begin{remark}[Внутреннее описание]
Голоморфное касательное пространство комплексного многообразия --- это пространство
дифференцирований $\cO^{\an}_{p}\to\C$ кольца ростков голоморфных функций, т.\,е.
$(\mm^{\an}_p/(\mm^{\an}_p)^2)^*$. Естественное отображение $\cO_{X,p}\to\cO^{\an}_{X,p}$ индуцирует
изоморфизмы $\cO_{X,p}/\mm_p^k\cong\cO^{\an}_{X,p}/(\mm^{\an}_p)^k$ при всех $k$ (в любой, даже
особой, точке; \cite{Se}), в частности $\mm_p/\mm_p^2\cong\mm_p^{\an}/(\mm_p^{\an})^2$. Поэтому
определение Зарисского --- это ровно «аналитическое касательное пространство Зарисского», которое
в гладких точках совпадает с касательным пространством многообразия.
\end{remark}

\subsection{Когда определения совпадают}

\begin{theorem}\label{thm:tangent-coincide}
Пусть $X$ --- многообразие размерности $d$, $p\in X$.
\begin{enumerate}
\item Если $p$ --- гладкая точка, то $X^{\an}$ около $p$ --- комплексное многообразие, и
$T^{\Zar}_pX=T_pX^{\an}$ (теорема~\ref{thm:smooth-manifold}).
\item Если $p$ --- особая точка, то $X^{\an}$ ни в какой окрестности $p$ не является комплексным
подмногообразием (и даже $C^1$-гладким вещественным подмногообразием размерности $2d$ в
$\C^n$), так что «дифференциально-геометрическое» касательное пространство не определено, а
$\dim T^{\Zar}_pX>d$.
\end{enumerate}
Итак, определения совпадают ровно на гладких многообразиях (в гладких точках).
\end{theorem}
\begin{proof}[Набросок доказательства (2)]
Если $X^{\an}$ около $p$ --- комплексное многообразие размерности $d$, то
$\cO^{\an}_{X,p}\cong\C\{z_1,\dots,z_d\}$ и $\dim\mm^{\an}_p/(\mm^\an_p)^2=d$; по предыдущему
замечанию $\dim\mm_p/\mm_p^2=d$, т.\,е. $p$ гладкая. Если же $X^{\an}$ --- лишь $C^1$-подмногообразие
с касательным пространством $V$, то $V$ совпадает с касательным конусом, который является
комплексным конусом; тогда $V$ --- комплексное подпространство размерности $d$, проекция
$X\to V$ --- локальный $C^1$-диффеоморфизм, и вместе с голоморфностью это даёт голоморфную карту
(подробности: теорема о том, что $C^1$-гладкое аналитическое множество голоморфно гладко, см.
\cite[гл.~1]{GH}).
\end{proof}

\begin{example}[Топологическое многообразие, не являющееся гладким]
Каспидальная кривая $X=\{y^2=x^3\}$ гомеоморфна $\C$ через $t\mapsto(t^2,t^3)$, т.\,е. $X^{\an}$
--- топологическое многообразие. Но $T^{\Zar}_0X=\C^2$. Она не является и $C^1$-подмногообразием:
иначе её касательная плоскость была бы (как предел секущих) прямой $\{y=0\}$, и проекция
$(x,y)\mapsto x$ была бы локальным диффеоморфизмом около $0$, а она двулистна:
$x=t^2$ имеет два прообраза $\pm t$.
\end{example}

\subsection{Касательное и кокасательное расслоения, дифференциальные формы}\label{ss:forms}

Пусть $X$ гладкое, $\dim X=d$. Функции $u_1,\dots,u_d\in\mm_p$, дифференциалы которых образуют
базис $\mm_p/\mm_p^2$, называются \emph{локальными параметрами} в $p$.

\begin{definition}[{\cite[II.3]{Sh}}]
\emph{Регулярная дифференциальная 1-форма} на открытом $U\subset X$ --- сопоставление $q\mapsto
\omega_q\in T_q^*X$, локально имеющее вид $\sum g_i\,df_i$ с регулярными $f_i,g_i$. Они
образуют пучок $\cO_X$-модулей $\Omega^1_X$; $\Omega^k_X=\Lambda^k\Omega^1_X$.
\end{definition}

\begin{proposition}\label{prop:local-params}
Если $u_1,\dots,u_d$ --- локальные параметры в $p$, то в некоторой окрестности $U\ni p$ формы
$du_1,\dots,du_d$ образуют базис $\Omega^1_X|_U$ над $\cO_U$. В частности, $\Omega^1_X$ локально
свободен ранга $d$, $\Omega^k_X$ локально свободен ранга $\binom dk$, а $\Omega^d_X$ обратим.
\end{proposition}
\begin{proof}
Пусть $X\subset\A^n$ аффинно, $u_i$ --- ограничения многочленов $U_i$, а
$f_1,\dots,f_{n-d}\in I(X)$ --- как в доказательстве теоремы~\ref{thm:smooth-manifold}. Тогда
дифференциалы $dU_1,\dots,dU_d,df_1,\dots,df_{n-d}$ линейно независимы в $p$ (образы $dU_i$ в
$(\C^n)^*/\langle d_pf_j\rangle=T_p^*X$ --- базис), т.\,е. определитель $\Delta$ их матрицы Якоби
по $x_1,\dots,x_n$ отличен от нуля в $p$. По правилу Крамера на $\{\Delta\neq0\}$
\[
dx_k=\sum_i a_{ki}\,dU_i+\sum_j b_{kj}\,df_j,\qquad a_{ki},b_{kj}\in\C[x][\Delta^{-1}].
\]
На $X$ имеем $df_j=0$ (так как $f_j|_X=0$), значит, $dx_k=\sum a_{ki}du_i$, и любая
$\sum g\,df=\sum g\,\frac{\partial f}{\partial x_k}dx_k$ выражается через $du_i$ с регулярными
коэффициентами на $U=X\cap\{\Delta\ne0\}$. Линейная независимость --- поточечная.
\end{proof}

\begin{definition}
\emph{Касательное расслоение} $T_X$ --- векторное расслоение, двойственное к $\Omega^1_X$ (его
слои --- $T_pX$); \emph{кокасательное} --- $T^*_X$ с пучком сечений $\Omega^1_X$. Если $u$ и $v$ ---
две системы локальных параметров, то $du_i=\sum_j\frac{\partial u_i}{\partial v_j}dv_j$ и
$\frac{\partial}{\partial v_j}=\sum_i\frac{\partial u_i}{\partial v_j}\frac{\partial}{\partial u_i}$; поэтому функции
перехода $T_X$ --- матрицы Якоби $(\partial u_i/\partial v_j)$, а функции перехода $T^*_X$ --- обратные
транспонированные к ним матрицы $(\partial v_j/\partial u_i)$; все они регулярны и обратимы.
Его аналитификация --- голоморфное (ко)касательное расслоение комплексного многообразия
$X^{\an}$.
\end{definition}

\emph{Рациональная $k$-форма} --- элемент $\Omega^k_X\otimes\C(X)$; около точки с локальными
параметрами $u$ она записывается как $\sum f_I\,du_{i_1}\wedge\dots\wedge du_{i_k}$ с
$f_I\in\C(X)$.

% ============================================================
\section{Канонический класс}\label{sec:canonical}
% ============================================================

Пусть $X$ --- гладкое многообразие размерности $n$.

\subsection{Три определения}

\begin{definition}
\begin{enumerate}
\item (\emph{Пучок.}) \emph{Канонический пучок} $\omega_X=\Omega^n_X=\det\Omega^1_X$ --- пучок
регулярных $n$-форм (обратим по предложению~\ref{prop:local-params}).
\item (\emph{Расслоение.}) \emph{Каноническое расслоение} $K_X=\det T^*_X=\Lambda^nT^*_X$. Если
$u$, $v$ --- локальные параметры на $U_i$, $U_j$, то
\[
du_1\wedge\dots\wedge du_n=\det\Big(\frac{\partial u}{\partial v}\Big)\,dv_1\wedge\dots\wedge dv_n,
\]
так что функции перехода $K_X$ --- якобианы $g_{ij}=\det(\partial v/\partial u)$
(в соглашении $s_i=g_{ij}s_j$ для коэффициентов формы $s_i\,du_1\wedge\ldots\wedge du_n$).
\item (\emph{Дивизор.}) Пусть $\omega\neq0$ --- рациональная $n$-форма. Около каждой точки
$\omega=f\,du_1\wedge\dots\wedge du_n$, $f\in\C(X)^*$. Положим
\[
\divv(\omega)=\sum_Y\ord_Y(f)\,Y\quad\text{($f$ берётся в карте, пересекающей $Y$)}.
\]
Это корректно: при смене параметров $f$ умножается на якобиан, регулярный и обратимый.
\emph{Канонический дивизор} --- любой $K_X=\divv(\omega)$; \emph{канонический класс} --- его
класс в $\Cl X=\Pic X$.
\end{enumerate}
\end{definition}

Рациональные $n$-формы существуют (например, $du_1\wedge\dots\wedge du_n$ для локальных
параметров на открытом $U$: это рациональная форма на $X$), и любые две отличаются на
рациональную функцию: $\omega'=h\omega$ (пространство рациональных $n$-форм одномерно над
$\C(X)$). Так как $\divv(h\omega)=\divv(h)+\divv(\omega)$, \emph{канонический класс определён
корректно}.

\begin{theorem}
Три определения согласованы: $\Omega^n_X$ --- пучок сечений расслоения $K_X$, и для любой
рациональной $n$-формы $\omega\ne0$
\[
\cO_X(\divv\omega)\xrightarrow{\ \sim\ }\Omega^n_X,\qquad h\mapsto h\,\omega.
\]
\end{theorem}
\begin{proof}
Первое утверждение --- определение: локальное сечение $\Lambda^nT^*_X$ над $U_i$ ---
это $s_i\,du^{(i)}_1\wedge\dots\wedge du^{(i)}_n$ с регулярным $s_i$. Для второго: на $U_i$
запишем $\omega=f_i\,\omega_i$, $\omega_i=du^{(i)}_1\wedge\dots\wedge du^{(i)}_n$ --- образующая
$\Omega^n|_{U_i}$. Тогда $\divv\omega$ --- дивизор Картье $\{(U_i,f_i)\}$ (отношения $f_i/f_j$ ---
якобианы), и $\cO(\divv\omega)|_{U_i}=f_i^{-1}\cO_{U_i}$. Для рациональной $h$:
$h\omega=(hf_i)\omega_i$ регулярна на $U_i$ $\iff$ $hf_i\in\cO(U_i)$ $\iff$
$h\in f_i^{-1}\cO(U_i)$. Это изоморфизм $\cO_{U_i}$-модулей, согласованный на пересечениях.
\end{proof}

Итак, $K_X$ можно понимать как класс дивизора, линейное расслоение или обратимый пучок
$\omega_X$ --- по теореме~\ref{thm:Weil=Cartier} это одно и то же. Для нормальных (особых)
многообразий канонический класс определяется как класс Вейля $\divv(\omega)$ на гладкой части
(дополнение к ней имеет коразмерность $\ge2$, и $\Cl X=\Cl(X\setminus\Sing X)$); он может не быть
классом Картье. Аналитически: $K_X$ --- голоморфное каноническое расслоение $\Lambda^nT^{*1,0}$, и
$c_1(X):=c_1(T_X)=-c_1(K_X)$.

\begin{proposition}\label{prop:K-product}
$K_{X\times Y}=p_1^*K_X\otimes p_2^*K_Y$.
\end{proposition}
\begin{proof}
$T^*_{X\times Y}=p_1^*T^*_X\oplus p_2^*T^*_Y$, и $\det(E\oplus F)=\det E\otimes\det F$
(лемма~\ref{lem:det}).
\end{proof}

\subsection{Проективное пространство}

\begin{theorem}\label{thm:KPn}
$K_{\PP^n}\cong\cO_{\PP^n}(-n-1)$, т.\,е. $K_{\PP^n}=-(n+1)H$.
\end{theorem}
\begin{proof}[Первое доказательство (через форму)]
В карте $U_0$ с координатами $x_i=X_i/X_0$ возьмём $\omega=dx_1\wedge\dots\wedge dx_n$; она
регулярна и нигде не обращается в нуль на $U_0$. Поэтому $\divv\omega$ сосредоточен на
$H_0=\{X_0=0\}$, и надо найти $\ord_{H_0}\omega$. В карте $U_1$ координаты
$y_1=X_0/X_1$, $y_j=X_j/X_1$ ($j\ge2$), $H_0\cap U_1=\{y_1=0\}$, и
\[
x_1=\frac1{y_1},\qquad x_j=\frac{y_j}{y_1}\ (j\ge2),\qquad
dx_1=-\frac{dy_1}{y_1^2},\qquad dx_j=\frac{dy_j}{y_1}-\frac{y_j\,dy_1}{y_1^2}.
\]
При внешнем умножении на $dx_1$ слагаемые с $dy_1$ в $dx_j$ пропадают, поэтому
\[
\omega=-\frac{dy_1}{y_1^2}\wedge\frac{dy_2}{y_1}\wedge\dots\wedge\frac{dy_n}{y_1}
=-\,y_1^{-(n+1)}\,dy_1\wedge\dots\wedge dy_n .
\]
Итак, $\divv\omega=-(n+1)H_0$.
\end{proof}

\begin{proof}[Второе доказательство (последовательность Эйлера)]
Докажем точность последовательности
\[
0\to\cO_{\PP^n}\xrightarrow{\ \epsilon\ }\cO_{\PP^n}(1)^{\oplus(n+1)}\xrightarrow{\ \beta\ } T_{\PP^n}\to0 .
\]
Пусть $\pi\colon\C^{n+1}\setminus0\to\PP^n$. Набор $(s_0,\dots,s_n)$ линейных (вообще,
однородных степени~1 по $X$) функций задаёт векторное поле $\sum s_i\partial/\partial X_i$ на
$\C^{n+1}\setminus 0$, инвариантное относительно гомотетий $X\mapsto\lambda X$ (коэффициенты
степени~1, а $\partial/\partial X_i$ степени $-1$); поэтому $d\pi$ переводит его в векторное поле
на $\PP^n$. Локально это определяет $\beta$ (для сечений $\cO(1)$ над $U\subset\PP^n$ ---
однородных рациональных функций степени 1 --- то же самое). Поле $d\pi(v)$ в точке $[X]$ равно нулю
ровно тогда, когда $v(X)\in\C X$, т.\,е. $v$ касается слоя $\pi$. Отсюда ядро $\beta$ порождено
полем Эйлера $\sum X_i\partial/\partial X_i$, т.\,е. образом $\epsilon(1)=(X_0,\dots,X_n)$.
Сюръективность: в карте $U_0$, $x_i=X_i/X_0$, имеем
$d\pi\big(X_0\frac{\partial}{\partial X_i}\big)=\frac{\partial}{\partial x_i}$, $i\ge1$.
Беря детерминанты (лемма~\ref{lem:det}):
$\cO(n+1)=\det\cO(1)^{n+1}\cong\det\cO\otimes\det T_{\PP^n}$, т.\,е. $\det T_{\PP^n}=\cO(n+1)$ и
$K_{\PP^n}=\cO(-n-1)$.
\end{proof}

\begin{example}
$K_{\PP^1}=\cO(-2)$ (на $\PP^1$ форма $dz$ имеет двойной полюс в $\infty$: $dz=-dw/w^2$);
$K_{\PP^1\times\PP^1}=\cO(-2,-2)$ по предложению~\ref{prop:K-product};
$K_{\PP^n\times\PP^m}=\cO(-n-1,-m-1)$.
\end{example}

\subsection{Грассманианы}

Пусть $V=\C^n$, $G=\Gr(k,n)$ --- многообразие $k$-мерных подпространств $W\subset V$,
$\dim G=k(n-k)$.

\emph{Карты.} Для подмножества $I\subset\{1,\dots,n\}$, $|I|=k$, пусть $U_I$ --- множество $W$,
проекция которых на $\C^I$ (вдоль $\C^{I^c}$) --- изоморфизм. Такое $W$ однозначно является
графиком линейного отображения $A\colon\C^I\to\C^{I^c}$, т.\,е. строками $k\times n$-матрицы с
единичной подматрицей в столбцах $I$. Так $U_I\cong\Hom(\C^I,\C^{I^c})\cong\A^{k(n-k)}$; функции
перехода рациональны (обратные миноры), что задаёт структуру гладкого многообразия.

\emph{Вложение Плюккера.} $W=\langle w_1,\dots,w_k\rangle\mapsto[w_1\wedge\dots\wedge w_k]\in
\PP(\Lambda^kV)$; координаты --- миноры $p_I$ матрицы со строками $w_i$. Образ замкнут (задаётся
соотношениями Плюккера), $U_I=\{p_I\ne0\}$, так что $G$ --- гладкое проективное многообразие.

\emph{Тавтологические расслоения.} $S\subset V\otimes\cO_G$ --- подрасслоение ранга $k$ со слоем
$W$ над $[W]$; $Q=(V\otimes\cO_G)/S$ --- \emph{факторрасслоение} ранга $n-k$:
\[
0\to S\to V\otimes\cO_G\to Q\to0.
\]
При вложении Плюккера $\Lambda^kS\subset\Lambda^kV\otimes\cO$ --- это ограничение
тавтологического расслоения $\cO_{\PP(\Lambda^kV)}(-1)$, т.\,е.
\[
\cO_G(1)\cong\det S^\vee,\qquad \det Q\cong(\det S)^{-1}\cong\cO_G(1)
\]
(второе --- из точной последовательности и леммы~\ref{lem:det}, так как $\det(V\otimes\cO)$
тривиально).

\begin{theorem}\label{thm:TGr}
$T_G\cong\mathcal{H}om(S,Q)=S^\vee\otimes Q$.
\end{theorem}
\begin{proof}
Зафиксируем $W_0\in G$ и дополнение $V=W_0\oplus W'$. Подпространства, трансверсальные $W'$,
--- это графики $\Gamma_A=\{w+Aw: w\in W_0\}$, $A\in\Hom(W_0,W')$; это аффинная карта (одна из
$U_I$ при координатных $W_0,W'$). Её дифференциал в $0$ даёт $T_{W_0}G\cong\Hom(W_0,W')
\cong\Hom(W_0,V/W_0)$. Проверим, что этот изоморфизм не зависит от выбора дополнения и
алгебраически зависит от точки. Опишем его внутренне: для кривой $W(t)$ с $W(0)=W$ и
$w\in W$ выберем кривую $w(t)\in W(t)$ с $w(0)=w$ и положим
\[
\varphi(w)=w'(0)\bmod W\ \in V/W.
\]
Корректность: если $\tilde w(t)$ --- другое поднятие, то $u(t)=\tilde w(t)-w(t)\in W(t)$,
$u(0)=0$, значит, $u(t)=t\,u_1(t)$ с $u_1(t)\in W(t)$ при $t\ne0$, а по непрерывности
$u'(0)=u_1(0)\in W(0)=W$. Линейность $\varphi$ по $w$ очевидна. В карте
$W(t)=\Gamma_{A(t)}$, $A(0)=A$, берём $w(t)=x+A(t)x$ ($x\in W_0$), и
$\varphi(x+Ax)=A'(0)x\bmod W$: соответствие $A'(0)\leftrightarrow\varphi$ --- линейный изоморфизм,
и в координатах карты это регулярная (линейная по $A'(0)$ с коэффициентами, полиномиальными по $A$)
замена. Значит, получаем изоморфизм векторных расслоений $T_G\cong\mathcal{H}om(S,Q)$.
\end{proof}

\begin{corollary}\label{cor:KGr}
$K_{\Gr(k,n)}\cong\cO_G(-n)$.
\end{corollary}
\begin{proof}
По лемме~\ref{lem:det}(2)
\[
\det T_G=\det(S^\vee\otimes Q)=(\det S^\vee)^{\otimes(n-k)}\otimes(\det Q)^{\otimes k}
=\cO_G(n-k)\otimes\cO_G(k)=\cO_G(n).\qedhere
\]
\end{proof}

\begin{example}
\begin{itemize}
\item $k=1$: $\Gr(1,n)=\PP^{n-1}$, $K=\cO(-n)$ --- согласуется с теоремой~\ref{thm:KPn}.
Последовательность Эйлера --- частный случай: $T=\mathcal Hom(S,Q)$, $S=\cO(-1)$, и
$0\to\mathcal Hom(S,S)=\cO\to\mathcal Hom(S,V\otimes\cO)=\cO(1)^n\to T\to0$.
\item $\Gr(2,4)\subset\PP^5$ --- квадрика Плюккера $p_{12}p_{34}-p_{13}p_{24}+p_{14}p_{23}=0$;
$K=\cO(-4)$. В~\S\ref{sec:adjunction} то же получится из формулы присоединения:
$K=\cO(-6+2)$.
\item Как дивизор: $K_G=-n\,\sigma_1$, где $\sigma_1=\{W:W\cap V_{n-k}\ne0\}$ --- дивизор Шуберта
(гиперплоское сечение $p_I=0$). Это можно проверить и прямо через форму
$\omega=\bigwedge da_{ij}$ на карте $U_I\cong\A^{k(n-k)}$: она не имеет нулей и полюсов на
$U_I$, а вдоль $G\setminus U_I=\{p_I=0\}$ имеет полюс порядка $n$.
\end{itemize}
\end{example}

\begin{remark}
Число $r$ такое, что $-K_X=rH$ с примитивным обильным $H$, называется \emph{индексом Фано}:
у $\PP^n$ индекс $n+1$, у $\Gr(k,n)$ --- $n$, у квадрики $Q^m\subset\PP^{m+1}$ ($m\ge2$) --- $m$
(см.~\S\ref{sec:adjunction}). Индекс Фано $n$-мерного многообразия не превосходит $n+1$.
\end{remark}

Другие примеры вычислений канонического класса --- для гиперповерхностей и полных пересечений,
торических многообразий и раздутий --- даны в \S\S\ref{sec:toric}, \ref{sec:adjunction},
\ref{sec:blowup}.

% ============================================================
\section{Отступление: торические многообразия}\label{sec:toric}
% ============================================================

Торические многообразия --- «полигон» алгебраической геометрии: все объекты предыдущих
разделов (дивизоры, расслоения, канонический класс, обильность, раздутия) вычисляются в них
комбинаторно по конусам и многогранникам. Основные источники: \cite{CLS,Fu,Da}.

% --- TikZ helpers
\newcommand{\lat}[1]{\foreach \x in {-#1,...,#1}\foreach \y in {-#1,...,#1}{\fill[gray!75] (\x,\y) circle (1.1pt);}}
\newcommand{\cfill}[5]{\fill[#5] (0,0) -- (10*#1,10*#2) -- (10*#3,10*#4) -- cycle;}
\newcommand{\ray}[2]{\draw[very thick,-{Stealth[length=2.2mm]}] (0,0) -- (#1,#2); \fill (#1,#2) circle (1.5pt);}

\subsection{Тор, характеры, однопараметрические подгруппы}

\emph{Алгебраический тор} $T=(\C^*)^n$ --- аффинное многообразие с координатным кольцом
$\C[t_1^{\pm1},\dots,t_n^{\pm1}]$ и групповой структурой. \emph{Характер} --- морфизм групп
$\chi\colon T\to\C^*$; \emph{однопараметрическая подгруппа} --- морфизм групп $\lambda\colon\C^*\to T$.

\begin{lemma}\label{lem:units-torus}
$\cO(T)^*=\{c\,t^m: c\in\C^*,\ m\in\Z^n\}$, где $t^m=t_1^{m_1}\cdots t_n^{m_n}$. Поэтому характеры
--- это мономы $\chi^m=t^m$, $m\in M:=\Z^n$, а однопараметрические подгруппы --- это
$\lambda^u(z)=(z^{u_1},\dots,z^{u_n})$, $u\in N:=\Z^n$.
\end{lemma}
\begin{proof}
Многогранник Ньютона произведения многочленов Лорана есть сумма Минковского многогранников
сомножителей. Если $fg=1$, то сумма многогранников Ньютона --- точка, значит, оба --- точки, т.\,е.
$f$ --- моном. Характер обратим и $\chi(1)=1$, поэтому $\chi=t^m$. Для $\lambda$ применяем то же к
$t_i\circ\lambda\in\cO(\C^*)^*$.
\end{proof}

Решётки $M$ (характеров) и $N$ (однопараметрических подгрупп) двойственны:
$\chi^m(\lambda^u(z))=z^{\langle m,u\rangle}$. Положим $M_\R=M\otimes\R$, $N_\R=N\otimes\R$.

\begin{definition}
\emph{Торическое многообразие} --- нормальное многообразие $X$, содержащее тор $T$ как плотное
открытое подмножество, такое что действие $T$ на себе сдвигами продолжается до морфизма
$T\times X\to X$.
\end{definition}

\subsection{Конусы и аффинные торические многообразия}

\begin{definition}
\emph{Рациональный полиэдральный конус} $\sigma=\Cone(u_1,\dots,u_s)=\{\sum\lambda_iu_i:\lambda_i\ge0\}
\subset N_\R$, $u_i\in N$; он \emph{строго выпуклый}, если $\sigma\cap(-\sigma)=\{0\}$.
\emph{Двойственный конус} $\sigma^\vee=\{m\in M_\R:\langle m,u\rangle\ge0\ \forall u\in\sigma\}$.
\emph{Грань} $\sigma$ --- $\tau=\sigma\cap m^\perp$ для $m\in\sigma^\vee$. \emph{Луч} --- одномерная
грань; его \emph{примитивный образующий} $u_\rho$ --- первая ненулевая точка решётки на луче.
\end{definition}

\begin{lemma}[Гордан]
Полугруппа $S_\sigma=\sigma^\vee\cap M$ конечно порождена.
\end{lemma}
\begin{proof}
$\sigma^\vee$ --- рациональный полиэдральный конус (двойственность выпуклых конусов):
$\sigma^\vee=\Cone(m_1,\dots,m_r)$, $m_i\in M$. Множество
$K=\{\sum\lambda_im_i:0\le\lambda_i<1\}$ ограничено, так что $K\cap M$ конечно. Любое
$m=\sum\lambda_im_i\in S_\sigma$ равно $\sum\lfloor\lambda_i\rfloor m_i+k$, $k\in K\cap M$.
Итак, $S_\sigma$ порождена $m_i$ и $K\cap M$.
\end{proof}

\begin{definition}
Пусть $\sigma$ строго выпуклый, $m_1,\dots,m_s$ --- образующие $S_\sigma$. \emph{Аффинное
торическое многообразие} $U_\sigma$ --- замыкание образа отображения
\[
\Phi\colon T\to\C^s,\qquad t\mapsto(\chi^{m_1}(t),\dots,\chi^{m_s}(t)).
\]
\end{definition}

\begin{proposition}
$\C[U_\sigma]\cong\C[S_\sigma]=\bigoplus_{m\in S_\sigma}\C\chi^m\subset\C[M]$; $\Phi$ ---
вложение $T$ как открытого плотного подмножества, и $U_\sigma$ --- торическое многообразие размерности
$n$. Тор действует по формуле $t\cdot(y_j)=(\chi^{m_j}(t)y_j)$.
\end{proposition}
\begin{proof}
Многочлен от $y_1,\dots,y_s$ обращается в нуль на $\Phi(T)$ (а значит, и на замыкании) тогда и
только тогда, когда его образ при $y_j\mapsto\chi^{m_j}$ равен нулю в $\C[M]$; поэтому
$\C[U_\sigma]$ --- образ этого гомоморфизма, т.\,е. $\C[S_\sigma]$. Так как $\sigma$ строго выпукл,
$\sigma^\vee$ полномерен, и $S_\sigma$ порождает $M$ как группу; значит, $\Phi$ инъективно, и
$\Phi(T)=\{y\in U_\sigma:\prod_j y_j\neq0\}$ открыто. Нормальность: $\C[S_\sigma]=\bigcap_\rho
\C[\{m:\langle m,u_\rho\rangle\ge0\}]$ --- пересечение нормальных колец \cite[1.3.5]{CLS}.
\end{proof}

\begin{example}\label{ex:toric-affine}
\begin{enumerate}
\item $\sigma=\Cone(e_1,\dots,e_k)\subset\R^n$: $S_\sigma$ порождена $e_1^*,\dots,e_k^*,
\pm e_{k+1}^*,\dots,\pm e_n^*$, и $U_\sigma=\C^k\times(\C^*)^{n-k}$. При $\sigma=\{0\}$
получаем $U_\sigma=T$.
\item $\sigma=\Cone(e_2,\,2e_1-e_2)\subset\R^2$ (рис.~\ref{fig:A1}). Тогда
$\sigma^\vee=\Cone((1,0),(1,2))$, базис Гильберта $S_\sigma$: $(1,0),(1,1),(1,2)$. Положив
$x=\chi^{(1,0)}$, $z=\chi^{(1,1)}$, $y=\chi^{(1,2)}$, получаем единственное соотношение
$xy=z^2$: $U_\sigma$ --- квадратичный конус из примера~\ref{ex:cone} (особенность $A_1$).
\end{enumerate}
\end{example}

\begin{figure}[ht]
\centering
\begin{tikzpicture}[scale=0.85]
\begin{scope}
\clip (-2.3,-2.3) rectangle (2.3,2.3);
\fill[blue!12] (0,0) -- (0,10) -- (20,-10) -- cycle;
\lat{2}
\draw[->,gray] (-2.3,0) -- (2.3,0); \draw[->,gray] (0,-2.3) -- (0,2.3);
\ray{0}{1}\ray{2}{-1}
\end{scope}
\node at (0.25,1.3) {$e_2$}; \node at (2.35,-1.3) {$2e_1-e_2$};
\node at (1.1,0.6) {$\sigma$};
\node at (0,-2.8) {$\sigma\subset N_\R$};
\end{tikzpicture}\qquad\qquad
\begin{tikzpicture}[scale=0.85]
\begin{scope}
\clip (-0.6,-0.6) rectangle (2.6,2.6);
\fill[red!12] (0,0) -- (3,0) -- (3*1,3*2) -- cycle;
\foreach \x in {0,...,3}\foreach \y in {0,...,3}{\fill[gray!75] (\x,\y) circle (1.1pt);}
\draw[->,gray] (-0.6,0) -- (2.6,0); \draw[->,gray] (0,-0.6) -- (0,2.6);
\draw[very thick] (0,0)--(3,0); \draw[very thick] (0,0)--(1.3,2.6);
\foreach \p in {(1,0),(1,1),(1,2)}{\fill[red!70!black] \p circle (2.2pt);}
\end{scope}
\node[right] at (1,-0.35) {$x$}; \node[right] at (1.05,1) {$z$}; \node[right] at (1.05,2) {$y$};
\node at (1,-1.3) {$\sigma^\vee\subset M_\R$,\ \ $xy=z^2$};
\end{tikzpicture}
\caption{Конус особенности $A_1$ и двойственный конус с базисом Гильберта $x,z,y$.}
\label{fig:A1}
\end{figure}

\subsection{Веера и общие торические многообразия}

Если $\tau=\sigma\cap m^\perp$, $m\in S_\sigma$, то $S_\tau=S_\sigma+\Z_{\ge0}(-m)$ и
$U_\tau=(U_\sigma)_{\chi^m}=\{\chi^m\ne0\}$ --- главное открытое подмножество $U_\sigma$.

\begin{definition}
\emph{Веер} $\Sigma$ в $N_\R$ --- конечный набор строго выпуклых рациональных полиэдральных
конусов, замкнутый относительно взятия граней и такой, что пересечение двух конусов --- грань каждого
из них. \emph{Торическое многообразие веера} $X_\Sigma$ получается склейкой $U_\sigma$,
$\sigma\in\Sigma$, по открытым подмножествам $U_\sigma\supset U_{\sigma\cap\sigma'}\subset
U_{\sigma'}$. Через $\Sigma(1)$ обозначаем множество лучей.
\end{definition}

\begin{theorem}[{\cite[3.1]{CLS}, \cite[1.4]{Fu}}]\label{thm:toric-basic}
\begin{enumerate}
\item Всякое (нормальное) торическое многообразие изоморфно $X_\Sigma$ для единственного
(с точностью до $\GL(N)$) веера $\Sigma$.
\item $X_\Sigma$ гладко $\iff$ каждый конус $\sigma\in\Sigma$ порождён частью базиса решётки $N$
(тогда $U_\sigma\cong\C^k\times(\C^*)^{n-k}$).
\item $X_\Sigma$ полно (т.\,е. $X^{\an}_\Sigma$ компактно) $\iff$ $|\Sigma|=\bigcup\sigma=N_\R$.
\item (\emph{Соответствие орбит и конусов.}) $T$-орбиты в $X_\Sigma$ взаимно однозначно
соответствуют конусам: $\sigma\leftrightarrow O(\sigma)\cong(\C^*)^{n-\dim\sigma}$, причём
$\overline{O(\sigma)}=\bigsqcup_{\tau\supseteq\sigma}O(\tau)$. В частности, лучам
$\rho\in\Sigma(1)$ отвечают $T$-инвариантные простые дивизоры $D_\rho=\overline{O(\rho)}$, а
$X_\Sigma\setminus T=\bigcup_\rho D_\rho$.
\end{enumerate}
\end{theorem}

Смысл веера хорошо виден через однопараметрические подгруппы:

\begin{proposition}
Для $u\in N$ предел $\lim_{z\to0}\lambda^u(z)$ существует в $U_\sigma$ тогда и только тогда, когда
$u\in\sigma$. Поэтому $X_\Sigma$ «знает» о направлениях $u$, вдоль которых тор можно продолжить.
\end{proposition}
\begin{proof}
Координаты точки $\lambda^u(z)$ в $U_\sigma\subset\C^s$ --- это $\chi^{m_j}(\lambda^u(z))=
z^{\langle m_j,u\rangle}$. Предел при $z\to0$ существует $\iff$ $\langle m_j,u\rangle\ge0$ при всех
$j$ $\iff$ $u\in(\sigma^\vee)^\vee=\sigma$.
\end{proof}

\begin{example}[Основные примеры, рис.~\ref{fig:fans}]\label{ex:fans}
\begin{enumerate}
\item $\PP^n$: лучи $e_1,\dots,e_n$, $e_0=-e_1-\dots-e_n$; максимальные конусы порождены любыми
$n$ из них. Конусу, не содержащему $e_i$, отвечает карта $U_i$, тор --- $\{x_0x_1\cdots x_n\ne0\}$,
$D_{e_i}=\{x_i=0\}$.
\item $\PP^1\times\PP^1$: лучи $\pm e_1,\pm e_2$ (произведение вееров).
\item Поверхность Хирцебруха $\mathbb F_a$: лучи $u_1=e_1$, $u_2=e_2$, $u_3=-e_1+ae_2$,
$u_4=-e_2$. Это $\PP^1$-расслоение $\PP(\cO\oplus\cO(a))$ над $\PP^1$; $\mathbb F_0=\PP^1\times\PP^1$,
а $\mathbb F_1=\Bl_p\PP^2$ (см.~\S\ref{sec:blowup}).
\item Раздутие $\PP^2$ в трёх торических неподвижных точках: шесть лучей
$\pm e_1,\pm e_2,\pm(e_1+e_2)$ («шестиугольник»).
\end{enumerate}
\end{example}

\begin{figure}[ht]
\centering
\begin{tikzpicture}[scale=0.62]
\begin{scope}
\clip (-2.2,-2.2) rectangle (2.2,2.2);
\cfill{1}{0}{0}{1}{blue!12}\cfill{0}{1}{-1}{-1}{green!12}\cfill{-1}{-1}{1}{0}{orange!15}
\lat{2}\ray{1}{0}\ray{0}{1}\ray{-1}{-1}
\end{scope}
\node at (0,-2.8) {$\PP^2$};
\end{tikzpicture}\hfill
\begin{tikzpicture}[scale=0.62]
\begin{scope}
\clip (-2.2,-2.2) rectangle (2.2,2.2);
\cfill{1}{0}{0}{1}{blue!12}\cfill{0}{1}{-1}{0}{green!12}\cfill{-1}{0}{0}{-1}{orange!15}\cfill{0}{-1}{1}{0}{red!10}
\lat{2}\ray{1}{0}\ray{0}{1}\ray{-1}{0}\ray{0}{-1}
\end{scope}
\node at (0,-2.8) {$\PP^1\times\PP^1$};
\end{tikzpicture}\hfill
\begin{tikzpicture}[scale=0.62]
\begin{scope}
\clip (-2.2,-2.2) rectangle (2.2,2.2);
\cfill{1}{0}{0}{1}{blue!12}\cfill{0}{1}{-1}{2}{green!12}\cfill{-1}{2}{0}{-1}{orange!15}\cfill{0}{-1}{1}{0}{red!10}
\lat{2}\ray{1}{0}\ray{0}{1}\ray{-1}{2}\ray{0}{-1}
\end{scope}
\node at (0,-2.8) {$\mathbb F_2$};
\end{tikzpicture}\hfill
\begin{tikzpicture}[scale=0.62]
\begin{scope}
\clip (-2.2,-2.2) rectangle (2.2,2.2);
\cfill{1}{0}{1}{1}{blue!12}\cfill{1}{1}{0}{1}{green!12}\cfill{0}{1}{-1}{0}{orange!15}
\cfill{-1}{0}{-1}{-1}{red!10}\cfill{-1}{-1}{0}{-1}{violet!12}\cfill{0}{-1}{1}{0}{yellow!20}
\lat{2}\ray{1}{0}\ray{1}{1}\ray{0}{1}\ray{-1}{0}\ray{-1}{-1}\ray{0}{-1}
\end{scope}
\node at (0,-2.8) {$\Bl_3\PP^2$};
\end{tikzpicture}
\caption{Веера $\PP^2$, $\PP^1\times\PP^1$, поверхности Хирцебруха $\mathbb F_2$ и раздутия $\PP^2$
в трёх точках. Цветные области --- двумерные конусы (карты $\C^2$), стрелки --- лучи
(инвариантные кривые $D_\rho$).}
\label{fig:fans}
\end{figure}

\subsection{Многогранники и проективные торические многообразия}

Пусть $P\subset M_\R$ --- полномерный решёточный многогранник. Для грани $F\subset P$ положим
$\sigma_F=\{u\in N_\R:\langle m,u\rangle=\min_{P}\langle\cdot,u\rangle\ \forall m\in F\}$ --- конус
внутренних нормалей. Конусы $\sigma_F$ образуют полный веер $\Sigma_P$ --- \emph{нормальный веер}
$P$; \emph{торическое многообразие многогранника} $X_P:=X_{\Sigma_P}$. При достаточно большом $k$
(для гладких $X_P$ --- уже при $k=1$, а для поверхностей всегда при $k=1$) отображение
\[
T\to\PP^{|kP\cap M|-1},\qquad t\mapsto[\chi^m(t)]_{m\in kP\cap M},
\]
продолжается до вложения $X_P$; в частности, $X_P$ проективно. Обратно, всякое проективное
торическое многообразие получается так \cite[2.3, 6.1]{CLS}.

\emph{Словарь:} вершины $P$ $\leftrightarrow$ максимальные конусы $\leftrightarrow$ неподвижные точки;
рёбра $\leftrightarrow$ $T$-инвариантные кривые; гиперграни $F_\rho=\{\langle m,u_\rho\rangle=-a_\rho\}$
$\leftrightarrow$ лучи $\rho$ $\leftrightarrow$ дивизоры $D_\rho$. Так, $\PP^2$ --- это треугольник, у
которого стороны --- три координатные прямые, а вершины --- три координатные точки.

\subsection{Дивизоры на торических многообразиях}

\begin{lemma}\label{lem:ord-toric}
Для луча $\rho$ и $m\in M$: $\ord_{D_\rho}(\chi^m)=\langle m,u_\rho\rangle$. Более общо, для многочлена
Лорана $f=\sum c_m\chi^m$ имеем $\ord_{D_\rho}(f)=\min\{\langle m,u_\rho\rangle: c_m\ne0\}$.
\end{lemma}
\begin{proof}
Дополним $u_\rho$ до базиса $e_1=u_\rho,e_2,\dots,e_n$ решётки $N$; пусть $e_i^*$ --- двойственный
базис, $x_i=\chi^{e_i^*}$. Тогда $S_\rho$ порождена $e_1^*,\pm e_2^*,\dots,\pm e_n^*$, и
$U_\rho=\C\times(\C^*)^{n-1}$ с координатами $x_1,x_2,\dots,x_n$, $D_\rho\cap U_\rho=\{x_1=0\}$.
Моном $\chi^m=x_1^{\langle m,e_1\rangle}x_2^{\langle m,e_2\rangle}\cdots$, а $x_2,\dots,x_n$ обратимы
на $U_\rho$; значит, его порядок вдоль $\{x_1=0\}$ равен $\langle m,u_\rho\rangle$. Для $f$ сгруппируем
$f=\sum_kx_1^kg_k(x_2,\dots,x_n)$; ненулевые $g_k$ --- многочлены Лорана от обратимых переменных,
не делящиеся на $x_1$; поэтому $\ord_{x_1=0}f$ --- наименьшее $k$ с $g_k\ne0$.
\end{proof}

\begin{theorem}\label{thm:Cl-toric}
Пусть лучи $u_\rho$ порождают $N_\R$ (например, $X_\Sigma$ полно). Тогда последовательность
\[
0\to M\xrightarrow{\ m\mapsto\divv(\chi^m)\ }\bigoplus_{\rho\in\Sigma(1)}\Z D_\rho\to\Cl X_\Sigma\to0,
\qquad \divv(\chi^m)=\sum_\rho\langle m,u_\rho\rangle D_\rho,
\]
точна. Если $X_\Sigma$ гладкое, то $\Pic X_\Sigma=\Cl X_\Sigma$; для гладкого полного $X_\Sigma$ это
свободная абелева группа ранга $|\Sigma(1)|-n$.
\end{theorem}
\begin{proof}
Формула для $\divv\chi^m$ --- лемма~\ref{lem:ord-toric} (у $\chi^m$ нет нулей и полюсов на $T$, а
$X\setminus T=\bigcup D_\rho$). Сюръективность: по предложению~\ref{prop:excision}
$\bigoplus\Z D_\rho\to\Cl X\to\Cl T\to0$ точна, а $\Cl T=0$, так как $\C[M]$ --- факториальное кольцо
(локализация кольца многочленов). Точность в среднем члене: если $\sum a_\rho D_\rho=\divv f$, то $f$
не имеет нулей и полюсов на $T$, т.\,е. $f\in\cO(T)^*$, и по лемме~\ref{lem:units-torus}
$f=c\chi^m$. Инъективность: если $\langle m,u_\rho\rangle=0$ для всех $\rho$, то $m=0$, так как $u_\rho$
порождают $N_\R$. Последнее --- теорема~\ref{thm:Weil=Cartier}; свобода $\Cl$ следует из того, что
в гладком полном случае некоторый максимальный конус порождён базисом $N$, и $D_\rho$ для лучей вне
этого конуса образуют базис $\Cl$.
\end{proof}

\begin{example}
\begin{itemize}
\item $\PP^n$: $\divv(\chi^{e_i^*})=D_i-D_0$; поэтому все $D_i$ линейно эквивалентны, и
$\Cl\PP^n=\Z$ (ср. пример~\ref{ex:ClPn}).
\item Конус $A_1$ (рис.~\ref{fig:A1}): $u_1=(0,1)$, $u_2=(2,-1)$; образ $M$ ---
$\{(m_2,\,2m_1-m_2)\}=\{(a,b): a+b\equiv0\bmod 2\}$, так что $\Cl U_\sigma=\Z^2/\text{образ}\cong
\Z/2$ (ср. пример~\ref{ex:cone}).
\item $\mathbb F_a$: соотношения $\divv\chi^{e_1^*}=D_1-D_3$, $\divv\chi^{e_2^*}=D_2+aD_3-D_4$;
$\Pic\mathbb F_a=\Z D_3\oplus\Z D_4$ ($D_3$ --- слой, $D_4$ --- сечение).
\end{itemize}
\end{example}

\begin{proposition}[{Дивизоры Картье, \cite[4.2]{CLS}}]
$T$-инвариантный дивизор $D=\sum a_\rho D_\rho$ является дивизором Картье тогда и только тогда, когда для
каждого конуса $\sigma\in\Sigma$ существует $m_\sigma\in M$ с $\langle m_\sigma,u_\rho\rangle=-a_\rho$ для
всех $\rho\in\sigma(1)$ (тогда на $U_\sigma$ $D=\divv(\chi^{-m_\sigma})$). В гладком случае это
выполнено всегда; для конуса $A_1$ дивизор $D_{u_1}$ не Картье ($2m_1=m_2=-1$ не имеет целых решений), а
$2D_{u_1}$ --- Картье.
\end{proposition}

\begin{theorem}[Глобальные сечения]\label{thm:toric-sections}
Для $D=\sum a_\rho D_\rho$ положим $P_D=\{m\in M_\R:\langle m,u_\rho\rangle\ge-a_\rho\ \forall\rho\}$. Тогда
\[
H^0(X_\Sigma,\cO(D))=\bigoplus_{m\in P_D\cap M}\C\cdot\chi^m .
\]
\end{theorem}
\begin{proof}
Глобальное сечение $\cO(D)$ --- это $f\in\C(X)$ с $\divv f+D\ge0$; в частности, $f$ регулярна на $T$,
т.\,е. $f=\sum c_m\chi^m$ --- многочлен Лорана. По лемме~\ref{lem:ord-toric} условие
$\ord_{D_\rho}f+a_\rho\ge0$ для всех $\rho$ равносильно тому, что $\langle m,u_\rho\rangle\ge-a_\rho$ для
всех $m$ с $c_m\ne0$ и всех $\rho$.
\end{proof}

\begin{example}
На $\PP^2$ для $D=dD_0$: $P_D=\{m_1\ge0,m_2\ge0,m_1+m_2\le d\}$ --- треугольник $d\Delta$, а
сечения --- мономы $x^{m_1}y^{m_2}$ степени $\le d$ в аффинной карте, т.\,е. однородные многочлены
степени $d$: $h^0(\cO(d))=\binom{d+2}2$.
\end{example}

\begin{proposition}[Пересечения на торических поверхностях]\label{prop:toric-surf}
Пусть $X_\Sigma$ --- гладкая полная торическая поверхность с лучами $u_1,\dots,u_r$ в циклическом
порядке, $D_i=D_{u_i}$. Тогда $u_{i-1}+u_{i+1}=b_iu_i$ для целого $b_i$, и
\[
D_i\cdot D_{i\pm1}=1,\qquad D_i^2=-b_i,\qquad D_i\cdot D_j=0\ \text{при}\ i-j\not\equiv0,\pm1\!\!\pmod r.
\]
\end{proposition}
\begin{proof}
Соседние $D_i,D_{i+1}$ пересекаются в неподвижной точке конуса $\Cone(u_i,u_{i+1})$; в карте
$U_\sigma\cong\C^2$ это координатные оси, пересекающиеся трансверсально: $D_i\cdot D_{i+1}=1$.
Несоседние дивизоры не пересекаются (нет конуса, содержащего оба луча). Для любого $m\in M$
\[
0=\divv(\chi^m)\cdot D_i=\langle m,u_{i-1}\rangle+\langle m,u_i\rangle D_i^2+\langle m,u_{i+1}\rangle,
\]
т.\,е. $\langle m,u_{i-1}+u_{i+1}+D_i^2u_i\rangle=0$ для всех $m$; значит, $u_{i-1}+u_{i+1}=-D_i^2u_i$.
\end{proof}

Например, у $\mathbb F_a$: $D_2^2=-a$, $D_4^2=a$, $D_1^2=D_3^2=0$; у шестиугольника все $D_i^2=-1$
(шесть $(-1)$-кривых, образующих цикл); у $\PP^2$ все $D_i^2=1$.

\begin{theorem}[{Обильность, \cite[6.1, 6.3]{CLS}}]\label{thm:toric-ample}
Пусть $X_\Sigma$ полно, $D=\sum a_\rho D_\rho$ --- дивизор Картье с данными $\{m_\sigma\}$.
\begin{enumerate}
\item $D$ обилен $\iff$ для каждого максимального конуса $\sigma$ и каждого $\rho\notin\sigma(1)$:
$\langle m_\sigma,u_\rho\rangle>-a_\rho$ («строгая выпуклость кусочно-линейной функции $D$»).
Тогда $P_D$ --- многогранник с нормальным веером $\Sigma$, а вершины $P_D$ --- это $m_\sigma$.
\item (\emph{Торический критерий Клеймана}.) Конус кривых порождён $T$-инвариантными кривыми; поэтому для
гладкого полного $X_\Sigma$: $D$ обилен $\iff$ $D\cdot C>0$ для всех $T$-инвариантных кривых $C$.
\end{enumerate}
\end{theorem}

\subsection{Канонический класс торического многообразия}

\begin{theorem}\label{thm:K-toric}
Для любого торического многообразия $X_\Sigma$:
\[
K_{X_\Sigma}=-\sum_{\rho\in\Sigma(1)}D_\rho .
\]
\end{theorem}
\begin{proof}
Рассмотрим $T$-инвариантную рациональную $n$-форму
\[
\Omega=\frac{dt_1}{t_1}\wedge\dots\wedge\frac{dt_n}{t_n}
\]
(здесь $t_i=\chi^{e_i^*}$ для базиса $e_i$ решётки $N$). Она регулярна и не обращается в нуль на $T$,
поэтому $\divv\Omega$ сосредоточен на $\bigcup D_\rho$, и достаточно найти $\ord_{D_\rho}\Omega$ (особые точки
имеют коразмерность $\ge2$ и на класс Вейля не влияют). Возьмём координаты $x_1,\dots,x_n$ на $U_\rho$
из леммы~\ref{lem:ord-toric}. Так как $t_j=\prod_ix_i^{a_{ij}}$ с унимодулярной матрицей
$A=(a_{ij})$ (переход между базисами $M$), имеем $\frac{dt_j}{t_j}=\sum_ia_{ij}\frac{dx_i}{x_i}$ и
\[
\Omega=\det A\cdot\frac{dx_1}{x_1}\wedge\dots\wedge\frac{dx_n}{x_n}
=\pm\frac{1}{x_1x_2\cdots x_n}\,dx_1\wedge\dots\wedge dx_n .
\]
На $U_\rho$ функции $x_2,\dots,x_n$ обратимы, а $x_1$ --- уравнение $D_\rho$: $\ord_{D_\rho}\Omega=-1$.
\end{proof}

\begin{corollary}
\begin{enumerate}
\item $K_{\PP^n}=-(D_0+\dots+D_n)\sim-(n+1)H$ --- третье доказательство теоремы~\ref{thm:KPn}.
\item $-K_{X_\Sigma}$ всегда эффективен; $K_{X_\Sigma}$ является дивизором Картье (многообразие
\emph{Горенштейново}) $\iff$ для каждого $\sigma$ есть $m_\sigma$ с $\langle m_\sigma,u_\rho\rangle=1$,
$\rho\in\sigma(1)$, т.\,е. образующие лучей каждого конуса лежат на аффинной гиперплоскости «высоты 1».
Для конуса $A_1$: $m=(1,1)$ подходит, так что $K$ --- Картье, хотя $D_{u_1}$ --- нет.
\item Для гладкой полной торической поверхности $-K\cdot D_i=2-b_i$ (из $K=-\sum D_j$ и
предложения~\ref{prop:toric-surf}; это же даёт формула присоединения: $K\cdot D_i=-2-D_i^2$).
\end{enumerate}
\end{corollary}

\begin{theorem}[Торические многообразия Фано]\label{thm:toric-Fano}
Гладкое полное торическое многообразие $X_\Sigma$ является многообразием Фано (т.\,е. $-K$ обилен, см.
\S\ref{sec:ample}) тогда и только тогда, когда точки $u_\rho$ являются вершинами выпуклого многогранника
$Q=\conv(u_\rho)$, а конусы $\Sigma$ --- это конусы над гранями $Q$. Тогда
$P_{-K}=\{m:\langle m,u_\rho\rangle\ge-1\}$ --- многогранник, полярный к $Q$, и
$h^0(-K)=|P_{-K}\cap M|$, $(-K)^n=n!\vol(P_{-K})$.
\end{theorem}
\begin{proof}
Для $D=-K$ все $a_\rho=1$. По теореме~\ref{thm:toric-ample} обильность означает, что для
каждого максимального $\sigma$ элемент $n_\sigma=-m_\sigma$ удовлетворяет $\langle n_\sigma,u_\rho\rangle=1$
при $\rho\in\sigma(1)$ и $\langle n_\sigma,u_\rho\rangle<1$ при $\rho\notin\sigma(1)$. Иначе говоря,
гиперплоскость $\{\langle n_\sigma,\cdot\rangle=1\}$ --- опорная к $Q$ и пересекает его ровно по
$\conv(u_\rho:\rho\in\sigma(1))$. Это и есть утверждение. Формула для объёма --- \cite[13.4.3]{CLS}.
\end{proof}

\begin{corollary}
Гладкая полная торическая поверхность --- поверхность дель Пеццо $\iff$ все $b_i\le1$. Таких поверхностей
ровно пять: $\PP^2$, $\PP^1\times\PP^1$, $\mathbb F_1=\Bl_1\PP^2$, $\Bl_2\PP^2$, $\Bl_3\PP^2$.
$\mathbb F_a$ при $a\ge2$ не Фано: $-K\cdot D_2=2-a\le0$.
\end{corollary}
\begin{proof}[Набросок]
Первое --- из теоремы~\ref{thm:toric-ample}(2) и $-K\cdot D_i=2-b_i$. Для гладкой полной торической
поверхности с $r$ лучами известно $\sum_i b_i=3r-12$ (это проверяется для $\PP^2$, $\mathbb F_a$ и
сохраняется при торических раздутиях, которыми получаются все такие поверхности \cite[10.4]{CLS}).
Из $b_i\le1$ следует $3r-12\le r$, т.\,е. $r\le6$; перебор даёт пять вееров рис.~\ref{fig:polytopes}.
\end{proof}

\begin{figure}[ht]
\centering
\newcommand{\polylat}{\foreach \x in {-1,...,2}\foreach \y in {-1,...,2}{\fill[gray!55] (\x,\y) circle (1pt);}}
\begin{tikzpicture}[scale=0.55]
\fill[blue!15,draw=blue!60!black,thick] (-1,-1)--(2,-1)--(-1,2)--cycle;
\polylat \fill (0,0) circle (2pt);
\node at (0.5,-2) {$\PP^2$, $K^2=9$};
\end{tikzpicture}\hfill
\begin{tikzpicture}[scale=0.55]
\fill[blue!15,draw=blue!60!black,thick] (-1,-1)--(1,-1)--(1,1)--(-1,1)--cycle;
\polylat \fill (0,0) circle (2pt);
\node at (0.5,-2) {$\PP^1\times\PP^1$, $8$};
\end{tikzpicture}\hfill
\begin{tikzpicture}[scale=0.55]
\fill[blue!15,draw=blue!60!black,thick] (0,-1)--(2,-1)--(-1,2)--(-1,0)--cycle;
\polylat \fill (0,0) circle (2pt);
\node at (0.5,-2) {$\Bl_1\PP^2$, $8$};
\end{tikzpicture}\hfill
\begin{tikzpicture}[scale=0.55]
\fill[blue!15,draw=blue!60!black,thick] (0,-1)--(1,-1)--(1,0)--(-1,2)--(-1,0)--cycle;
\polylat \fill (0,0) circle (2pt);
\node at (0.5,-2) {$\Bl_2\PP^2$, $7$};
\end{tikzpicture}\hfill
\begin{tikzpicture}[scale=0.55]
\fill[blue!15,draw=blue!60!black,thick] (1,-1)--(1,0)--(0,1)--(-1,1)--(-1,0)--(0,-1)--cycle;
\polylat \fill (0,0) circle (2pt);
\node at (0.5,-2) {$\Bl_3\PP^2$, $6$};
\end{tikzpicture}
\caption{Антиканонические многоугольники $P_{-K}$ пяти гладких торических поверхностей дель Пеццо.
Число граничных точек решётки равно $K^2$ (по формуле Пика площадь равна $K^2/2$), число всех точек ---
$h^0(-K)=K^2+1$. Каждый срезанный угол --- раздутие неподвижной точки.}
\label{fig:polytopes}
\end{figure}

\begin{remark}
\begin{itemize}
\item Многогранники $P$, у которых $P$ и полярный $P^*$ --- решёточные, называются \emph{рефлексивными}; они
отвечают торическим Фано с горенштейновыми особенностями. В размерности 2 их 16, и для рефлексивного
многоугольника число граничных точек $P$ плюс число граничных точек $P^*$ равно $12$.
\item Торическое раздутие неподвижной точки, отвечающей конусу $\Cone(u_i,u_{i+1})$, --- это
\emph{звёздное подразбиение}: добавляется луч $u_i+u_{i+1}$ (см. \S\ref{sec:blowup}, рис.~\ref{fig:star}).
На многоугольнике это срезание вершины.
\end{itemize}
\end{remark}

% ============================================================
\section{Формула присоединения}\label{sec:adjunction}
% ============================================================

\subsection{Формулировка и доказательство}

\begin{theorem}[Формула присоединения]\label{thm:adjunction}
Пусть $X$ --- гладкое многообразие, $Y\subset X$ --- гладкий простой дивизор. Тогда
\[
\omega_Y\cong\big(\omega_X\otimes\cO_X(Y)\big)\big|_Y,\qquad\text{т.\,е.}\qquad K_Y=(K_X+Y)|_Y .
\]
\end{theorem}

Доказательство состоит из двух лемм.

\begin{lemma}[Нормальное расслоение]\label{lem:normal}
Пусть $N_{Y/X}=T_X|_Y/T_Y$ --- нормальное расслоение. Тогда $N_{Y/X}\cong\cO_X(Y)|_Y$.
\end{lemma}
\begin{proof}
По теореме~\ref{thm:Weil=Cartier} $Y$ --- дивизор Картье: есть покрытие $X=\bigcup U_i$ и
$f_i\in\cO(U_i)$, $\divv(f_i)=Y\cap U_i$, причём (доказательство той же теоремы) $f_i$ порождает идеал
функций, равных нулю на $Y\cap U_i$; на пересечениях $f_i=g_{ij}f_j$, $g_{ij}\in\cO^*(U_i\cap U_j)$
--- функции перехода $\cO_X(Y)$ (лемма~\ref{lem:OD}).

Пусть $p\in Y\cap U_i$. По теореме~\ref{thm:jacobian} (применённой в аффинной окрестности, где идеал $Y$
порождён $f_i$ и уравнениями $X$) $T_pY=\{v\in T_pX: d_pf_i(v)=0\}$. Так как $\dim T_pY=\dim X-1$,
имеем $d_pf_i\neq0$. Значит, отображение расслоений
\[
\alpha_i\colon T_X|_{Y\cap U_i}\to\cO_{Y\cap U_i},\qquad v\mapsto df_i(v),
\]
сюръективно и имеет ядро $T_Y|_{U_i}$. На пересечениях
$df_i=g_{ij}\,df_j+f_j\,dg_{ij}=g_{ij}\,df_j$ на $Y$ (там $f_j=0$), т.\,е. $\alpha_i=g_{ij}\alpha_j$.
Это ровно означает, что $\{\alpha_i\}$ склеиваются в сюръективный морфизм расслоений
$\alpha\colon T_X|_Y\to\cO_X(Y)|_Y$ с ядром $T_Y$. Итак, $N_{Y/X}\cong\cO_X(Y)|_Y$.
\end{proof}

\begin{proof}[Доказательство теоремы~\ref{thm:adjunction}]
Из точной последовательности $0\to T_Y\to T_X|_Y\to N_{Y/X}\to0$ и леммы~\ref{lem:det}:
$\det T_X|_Y\cong\det T_Y\otimes N_{Y/X}$. Переходя к двойственным,
$\omega_X|_Y\cong\omega_Y\otimes N_{Y/X}^\vee=\omega_Y\otimes\cO_X(-Y)|_Y$, что и требовалось.
\end{proof}

Двойственная форма той же последовательности --- \emph{конормальная последовательность}
$0\to\cI_Y/\cI_Y^2\to\Omega^1_X|_Y\to\Omega^1_Y\to0$ с $\cI_Y/\cI_Y^2\cong\cO_X(-Y)|_Y$
(класс $f_i$ --- образующая).

\begin{remark}[Вычет Пуанкаре: явный изоморфизм]
Изоморфизм можно записать явно. Пусть около точки $Y=\{f=0\}$, $z_1,\dots,z_n$ --- локальные параметры
на $X$, и $\partial f/\partial z_n\ne0$ (так что $z_1,\dots,z_{n-1}$ --- параметры на $Y$). Для
формы $\eta=\dfrac{g\,dz_1\wedge\dots\wedge dz_n}{f}$ (сечения $\omega_X(Y)$) положим
\[
\Res_Y\eta=(-1)^{n-1}\frac{g\,dz_1\wedge\dots\wedge dz_{n-1}}{\partial f/\partial z_n}\Big|_Y .
\]
Это не зависит от выбора (проверка: $\eta=\frac{df}{f}\wedge\eta'$, $\Res\eta=\eta'|_Y$), и
$\Res\colon\omega_X(Y)|_Y\to\omega_Y$ --- изоморфизм. Например, для гладкой плоской кривой
$C=\{f(x,y)=0\}$ степени $d$ регулярные 1-формы на $C$ --- это
\[
\omega=\frac{P(x,y)\,dx}{\partial f/\partial y},\qquad \deg P\le d-3,
\]
откуда $g(C)=\dim H^0(\omega_C)=\binom{d-1}2$.
\end{remark}

\begin{corollary}\label{cor:adj-general}
\begin{enumerate}
\item Если $Y\subset X$ --- гладкое полное пересечение гладких дивизоров $D_1,\dots,D_r$
(трансверсально пересекающихся), то $K_Y=(K_X+D_1+\dots+D_r)|_Y$.
\item Если $Y\subset X$ --- гладкое подмногообразие любой коразмерности, то
$\omega_Y\cong\omega_X|_Y\otimes\det N_{Y/X}$.
\end{enumerate}
\end{corollary}
\begin{proof}
(1) --- индукция по $r$. (2) --- то же рассуждение с детерминантами.
\end{proof}

\subsection{Применения}

\begin{example}[Гиперповерхности в $\PP^n$]\label{ex:hypersurface}
Для гладкой гиперповерхности $X_d\subset\PP^n$ степени $d$ имеем $\cO(X_d)=\cO(d)$ и
\[
K_{X_d}=\cO_{X_d}(d-n-1).
\]
Отсюда трихотомия:
\begin{itemize}
\item $d\le n$: $-K$ обилен --- \emph{многообразие Фано} (квадрики, кубические поверхности,
кубические и квартичные трёхмерные многообразия\ldots);
\item $d=n+1$: $K\cong\cO$ --- \emph{многообразие Калаби--Яу}: плоская кубика (эллиптическая кривая),
квартика в $\PP^3$ (поверхность K3), квинтика в $\PP^4$;
\item $d\ge n+2$: $K$ обилен --- многообразие общего типа.
\end{itemize}
В частности, для гладкой \emph{кубической поверхности} $S\subset\PP^3$: $K_S=\cO_S(-1)$, т.\,е.
$-K_S$ --- гиперплоское сечение; $K_S^2=\deg S=3$.
\end{example}

\begin{example}[Род плоской кривой]
Для гладкой плоской кривой $C$ степени $d$: $K_C=\cO_C(d-3)$, $\deg\cO_C(1)=d$ (по Безу), поэтому
$2g-2=d(d-3)$ и
\[
g=\frac{(d-1)(d-2)}2 .
\]
Прямые и коники рациональны, гладкие кубики --- эллиптические кривые ($g=1$), гладкие квартики имеют род 3.
\end{example}

\begin{example}[Полные пересечения]
Для гладкого полного пересечения $X=X_{d_1,\dots,d_r}\subset\PP^n$: $K_X=\cO_X(\sum d_i-n-1)$.
Пересечение двух квадрик в $\PP^4$ --- поверхность с $K=\cO(-1)$, т.\,е. поверхность дель Пеццо
степени~4. Грассманиан $\Gr(2,4)$ --- квадрика в $\PP^5$: $K=\cO(2-6)=\cO(-4)$, что согласуется со
следствием~\ref{cor:KGr}. Квадрика $Q^m\subset\PP^{m+1}$ имеет $K=\cO(-m)$; в частности, для
$Q^2\cong\PP^1\times\PP^1$ (вложенной по Сегре, $\cO_Q(1)=\cO(1,1)$) получаем $K=\cO(-2,-2)$.
\end{example}

\begin{theorem}[Формула рода на поверхности]\label{thm:genus-formula}
Пусть $C$ --- гладкая кривая на гладкой проективной поверхности $S$. Тогда
\[
2g(C)-2=C\cdot(C+K_S).
\]
Для произвольной неприводимой кривой $C$ число $p_a(C)=1+\frac12C\cdot(C+K_S)$ (\emph{арифметический
род}) --- целое неотрицательное, и $p_a(C)=0$ тогда и только тогда, когда $C$ --- гладкая рациональная
кривая \cite[V, упр.~1.3]{Ha}.
\end{theorem}
\begin{proof}
$2g-2=\deg K_C=\deg(K_S+C)|_C=(K_S+C)\cdot C$ по теореме~\ref{thm:adjunction} и
теореме~\ref{thm:intersection}(2).
\end{proof}

\begin{example}\label{ex:adj-surfaces}
\begin{enumerate}
\item На $\PP^1\times\PP^1$ гладкая кривая бистепени $(a,b)$: $C^2=2ab$, $K\cdot C=-2a-2b$, откуда
$g=(a-1)(b-1)$.
\item \emph{$(-1)$-кривые.} Если $E\cong\PP^1$ и $E^2=-1$, то $K_S\cdot E=-1$. Обратно, неприводимая кривая
с $E^2<0$ и $K_S\cdot E<0$ --- это $(-1)$-кривая: $E^2+K\cdot E=2p_a-2\ge-2$ и обе величины $\le-1$.
\item \emph{Прямая на гладкой поверхности степени $d$ в $\PP^3$.} $K_S=(d-4)H|_S$, $H\cdot\ell=1$, $g=0$:
$-2=\ell^2+d-4$, т.\,е. $\ell^2=2-d$. При $d=2$: $\ell^2=0$ (образующие квадрики), при $d=3$:
$\ell^2=-1$ (\emph{прямые на кубике --- $(-1)$-кривые}), при $d=4$: $\ell^2=-2$.
\item \emph{Кривые на кубической поверхности.} $K_S=-H$, поэтому для неприводимой кривой
$C\subset S$: $\deg C=H\cdot C=-K_S\cdot C$. Кривая с $-K_S\cdot C=1$ --- прямая.
\end{enumerate}
\end{example}

% ============================================================
\section{Обильность. Поверхности дель Пеццо и многообразия Фано}\label{sec:ample}
% ============================================================

\subsection{Алгебраическая обильность}

\begin{definition}
Пусть $X$ --- проективное многообразие. Линейное расслоение $L$ \emph{очень обильно}, если существует
замкнутое вложение $\varphi\colon X\hookrightarrow\PP^N$ с $\varphi^*\cO(1)\cong L$ (равносильно: некоторые
сечения $s_0,\dots,s_N\in H^0(L)$ без общих нулей задают по \S\ref{ss:linsys} вложение). $L$
\emph{обильно}, если $L^{\otimes m}$ очень обильно для некоторого $m>0$. Дивизор $D$ (очень) обилен, если
таков $\cO(D)$.
\end{definition}

Вложение задаётся сечениями, если они \emph{разделяют точки} (для $p\ne q$ есть сечение, равное нулю в $p$
и не равное в $q$) и \emph{разделяют касательные векторы} (сечения, равные нулю в $p$, имеют
дифференциалы, порождающие $T_p^*X$) \cite[II.7.3]{Ha}.

\begin{theorem}[Основные критерии обильности]\label{thm:ample-criteria}
Пусть $X$ --- проективное многообразие, $L$ --- линейное расслоение.
\begin{enumerate}
\item (\emph{Серр--Картан--Гротендик}, \cite[II.7.6, III.7.5]{Ha}.) $L$ обильно $\iff$ для любого
когерентного пучка $\cF$ при $m\gg0$ пучок $\cF\otimes L^m$ порождён глобальными сечениями
$\iff$ $H^i(X,\cF\otimes L^m)=0$ при $i>0$, $m\gg0$.
\item (\emph{Накаи--Мойшезон}, \cite[V.1.10]{Ha}, \cite[1.2.23]{La}.) $L$ обильно $\iff$
$(L^{\dim V}\cdot V)>0$ для любого замкнутого подмногообразия $V\subset X$ положительной размерности.
Для поверхности: $L^2>0$ и $L\cdot C>0$ для всех неприводимых кривых $C$.
\item (\emph{Клейман}, \cite[1.4.23]{La}.) $L$ обильно $\iff$ $L\cdot z>0$ для всех ненулевых $z$ из
замкнутого конуса кривых $\NE(X)\subset N_1(X)_\R$.
\item На гладкой кривой: $L$ обильно $\iff\deg L>0$; при $\deg L\ge2g+1$ оно очень обильно.
\end{enumerate}
Кроме того: ограничение обильного на подмногообразие обильно; обратный образ обильного при конечном
морфизме обилен; тензорное произведение обильного и порождённого сечениями обильно.
\end{theorem}

Направление «$\Rightarrow$» в критерии Накаи--Мойшезона очевидно: если $L^m=\varphi^*\cO(1)$, то
$m^{\dim V}(L^{\dim V}\cdot V)=\deg\varphi(V)>0$.

\begin{example}\label{ex:ample-examples}
\begin{enumerate}
\item $\cO_{\PP^n}(d)$ обильно $\iff d>0$ ($v_d$ --- вложение).
\item $\cO(a,b)$ на $\PP^1\times\PP^1$ обильно $\iff a,b>0$. Действительно, $\cO(a,b)\cdot F_1=b$,
$\cO(a,b)\cdot F_2=a$ для слоёв $F_1,F_2$ проекций; если $a,b>0$, то $\cO(a,b)$ --- обратный образ $\cO(1)$
при композиции $v_a\times v_b$ и вложения Сегре.
\item На $X=\Bl_p\PP^2$ (\S\ref{sec:blowup}) с $\Pic X=\Z L\oplus\Z E$: $aL-bE$ обилен $\iff a>b>0$.
Неприводимая кривая $C\ne E$ имеет класс $dL-mE$ с $d\ge m\ge0$ ($m$ --- кратность в $p$ её образа
степени $d$), и $dL-mE=m(L-E)+(d-m)L$, $L=(L-E)+E$; поэтому конус кривых порождён $E$ и $L-E$
(собственный прообраз прямой через $p$). Имеем $(aL-bE)\cdot E=b$, $(aL-bE)\cdot(L-E)=a-b$, и критерий
Клеймана даёт ответ.
\end{enumerate}
\end{example}

\subsection{Положительность в дифференциальной геометрии}

Пусть $X$ --- компактное комплексное многообразие, $L$ --- голоморфное линейное расслоение с
эрмитовой метрикой $h$. Для локального голоморфного ненулевого сечения $e$ положим $|e|_h^2=h(e,e)$.
\emph{Кривизна} связности Черна
\[
\Theta_h=-\partial\bar\partial\log|e|_h^2
\]
не зависит от $e$ (замена $e\mapsto ge$ добавляет $\log|g|^2$, а $\partial\bar\partial\log|g|^2=0$ для
голоморфной $g\ne0$). Форма $c_1(L,h)=\frac{i}{2\pi}\Theta_h$ --- замкнутая вещественная $(1,1)$-форма,
представляющая $c_1(L)$ в $H^2(X,\R)$; для кривой $C\subset X$: $\deg L|_C=\int_Cc_1(L,h)$.

\begin{definition}
$L$ \emph{положительно}, если существует метрика $h$, для которой $\frac{i}{2\pi}\Theta_h$ --- положительно
определённая $(1,1)$-форма, т.\,е. кэлерова форма. В локальных координатах это означает, что матрица
$\big(\partial^2\varphi/\partial z_j\partial\bar z_k\big)$ положительно определена, где
$\varphi=-\log|e|^2_h$ (функция $\varphi$ строго плюрисубгармонична).
\end{definition}

\begin{example}[Метрика Фубини--Штуди]\label{ex:FS}
Тавтологическое расслоение $\cO(-1)\subset\PP^n\times\C^{n+1}$ наследует стандартную метрику $\C^{n+1}$.
В карте $U_0$ сечение $e(z)=(1,z_1,\dots,z_n)$, $|e|^2=1+|z|^2$, и
\[
\tfrac{i}{2\pi}\Theta_{\cO(-1)}=-\tfrac{i}{2\pi}\partial\bar\partial\log(1+|z|^2)=-\omega_{FS},\qquad
\omega_{FS}:=\tfrac{i}{2\pi}\partial\bar\partial\log(1+|z|^2).
\]
Двойственная метрика на $\cO(1)$ имеет кривизну $+\omega_{FS}$. Форма $\omega_{FS}$ положительна: в точке
$z=0$ имеем $\partial\bar\partial\log(1+|z|^2)=\sum dz_j\wedge d\bar z_j$, а $\frac{i}{2\pi}\sum dz_j\wedge
d\bar z_j$ положительна; в остальных точках --- по однородности ($U(n+1)$ действует транзитивно и
сохраняет $\omega_{FS}$). Итак, $\cO(1)$ положительно. Нормировка: при $n=1$
\[
\int_{\PP^1}\omega_{FS}=\frac{i}{2\pi}\int_\C\frac{dz\wedge d\bar z}{(1+|z|^2)^2}
=\frac1\pi\int_0^{2\pi}\!\!\int_0^\infty\frac{r\,dr\,d\theta}{(1+r^2)^2}=1=\deg\cO(1).
\]
\end{example}

\begin{theorem}[{Кодаира, теорема о вложении; \cite[гл.~1, \S\,4]{GH}, \cite[5.3]{Hu}}]\label{thm:Kodaira}
Для компактного комплексного многообразия $X$ и голоморфного линейного расслоения $L$:
$L$ положительно $\iff$ $L^{\otimes m}$ задаёт голоморфное вложение $X\hookrightarrow\PP^N$ при $m\gg0$.
В частности, $X$ проективно (по Чжоу) и, по GAGA, $L$ алгебраично и обильно. Итак, для гладкого
проективного $X$:
\[
\text{обильность (алгебраическое понятие)}\iff\text{положительность (дифференциально-геометрическое понятие)}.
\]
\end{theorem}
Направление «$\Leftarrow$»: если $L^m=\varphi^*\cO(1)$, то метрика $\varphi^*h_{FS}^{1/m}$ имеет кривизну
$\frac1m\varphi^*\omega_{FS}>0$ (так как $\varphi$ --- погружение). Направление «$\Rightarrow$» использует
теорему Кодаиры об обращении в нуль $H^q(X,K_X\otimes L)=0$, $q>0$, для положительного $L$.

\begin{remark}[Кривизна Риччи и антиканоническое расслоение]
Кэлерова метрика $\omega=\frac i2\sum g_{j\bar k}dz_j\wedge d\bar z_k$ задаёт метрику на
$K_X^{-1}=\Lambda^nT_X$: $|\partial_1\wedge\dots\wedge\partial_n|^2=\det(g_{j\bar k})$. Её кривизна ---
\emph{форма Риччи}
\[
\Ric(\omega)=-i\,\partial\bar\partial\log\det(g_{j\bar k}),\qquad [\Ric\omega]=2\pi c_1(X).
\]
Следовательно, если у $X$ есть кэлерова метрика с положительной кривизной Риччи, то $-K_X$ положительно,
т.\,е. обильно. Обратно, по теореме Яу (гипотеза Калаби) для любого Фано $X$ и любой кэлеровой формы в
классе $2\pi c_1(X)$ найдётся кэлерова метрика $\omega'$ с $\Ric(\omega')=$ этой форме $>0$. Пример:
$\Ric(\omega_{FS})=2\pi(n+1)\omega_{FS}$ на $\PP^n$ (метрика Кэлера--Эйнштейна).
\end{remark}

\subsection{Многообразия Фано}

\begin{definition}
Гладкое проективное многообразие $X$ называется \emph{многообразием Фано}, если антиканонический класс
$-K_X$ обилен (эквивалентно: $X$ допускает кэлерову метрику с $\Ric>0$). Двумерные многообразия Фано
называются \emph{поверхностями дель Пеццо}; их \emph{степень} --- $d=K_X^2$.
\end{definition}

\begin{theorem}[Свойства многообразий Фано]
Пусть $X$ --- многообразие Фано размерности $n$. Тогда
\begin{enumerate}
\item $H^i(X,\cO_X)=0$ при $i>0$ (по Кодаире: $\cO_X=K_X\otimes(-K_X)$), поэтому
$\Pic X\cong H^2(X,\Z)$ --- конечно порождённая группа (экспоненциальная последовательность);
\item $X$ односвязно и рационально связно (Кампана; Коллар--Миока--Мори); в размерности $\le2$ ---
рационально;
\item при каждом $n$ существует лишь конечное число деформационных семейств многообразий Фано;
при $n=1$ это $\PP^1$, при $n=2$ --- 10 семейств поверхностей дель Пеццо, при $n=3$ --- 105 семейств
(Исковских, Мори--Мукаи) \cite{IP}.
\end{enumerate}
\end{theorem}

\begin{theorem}[{Классификация поверхностей дель Пеццо; \cite[IV]{Ma}, \cite{De}, \cite[8]{Do}}]\label{thm:dP}
Поверхность дель Пеццо $S$ степени $d$ удовлетворяет $1\le d\le9$ и изоморфна либо
$\PP^1\times\PP^1$ ($d=8$), либо раздутию $\PP^2$ в $r=9-d\le8$ точках \emph{общего положения}:
никакие 3 не лежат на прямой, никакие 6 --- на конике, и (при $r=8$) не существует кубики, проходящей через
все 8 точек и особой в одной из них. Обратно, все эти поверхности --- дель Пеццо.
\end{theorem}

Для $r\le6$ мы докажем «обратно» в \S\ref{sec:blowup} (теорема~\ref{thm:dP-blowup}); там же объясняется,
почему условия общего положения необходимы. Схема доказательства прямого утверждения: (i) $S$ рациональна
по критерию Кастельнуово ($q=P_2=0$, так как $-K$ обилен); (ii) минимальная рациональная поверхность ---
это $\PP^2$ или $\mathbb F_a$ ($a\ne1$), а из них Фано только $\PP^2$ и $\mathbb F_0=\PP^1\times\PP^1$;
(iii) значит, $S$ получается раздутиями из $\PP^2$ или $\PP^1\times\PP^1$, и при нарушении общего положения
появляется кривая $C$ с $-K\cdot C\le0$.

\begin{remark}[Антиканонические модели]
При $d\ge3$ класс $-K_S$ очень обилен и вкладывает $S$ в $\PP^d$ как поверхность степени $d$: $d=3$ ---
кубическая поверхность в $\PP^3$, $d=4$ --- пересечение двух квадрик в $\PP^4$, $d=9$ --- $v_3(\PP^2)\subset\PP^9$.
При $d=2$ отображение $|-K|$ --- двулистное накрытие $\PP^2$, разветвлённое в квартике; при $d=1$ ---
$|-2K|$ даёт двулистное накрытие квадратичного конуса, разветвлённое в кривой степени 6 (не проходящей через вершину) и в вершине конуса.
\end{remark}

\subsection{Как доказать, что многообразие --- Фано}

Перечислим основные методы (с примерами).

\begin{enumerate}
\item \textbf{Явное вычисление $K_X$ через формулу присоединения.} Гладкая гиперповерхность
$X_d\subset\PP^n$ с $d\le n$: $-K=\cO_X(n+1-d)$ --- ограничение обильного (пример~\ref{ex:hypersurface}).
Полные пересечения с $\sum d_i\le n$. Гладкие гиперповерхности во взвешенных проективных
пространствах (например, гладкая поверхность степени 4 в $\PP(1,1,1,2)$ --- дель Пеццо степени 2, степени 6
в $\PP(1,1,2,3)$ --- степени 1).
\item \textbf{Однородные многообразия.} $\Gr(k,n)$: $-K=\cO(n)$, а $\cO(1)$ очень обилен (Плюккер). Вообще,
все флаговые многообразия $G/P$ --- Фано: $T_X$ порождено глобальными сечениями (векторными полями из
алгебры Ли $G$), поэтому и $-K=\det T_X$ порождено сечениями, а положительность проверяется на
$T$-инвариантных кривых.
\item \textbf{Произведения}: $X\times Y$ Фано $\iff$ $X$ и $Y$ Фано
($-K_{X\times Y}=p_1^*(-K_X)\otimes p_2^*(-K_Y)$).
\item \textbf{Численные критерии (Накаи--Мойшезон, Клейман).} Для поверхностей: проверить $K^2>0$ и
$-K\cdot C>0$ для всех неприводимых кривых. Обычно достаточно знать образующие конуса кривых.
Например, для $\Bl_p\PP^2$: $-K=3L-E$, $(-K)^2=8$, $-K\cdot E=1$, $-K\cdot(L-E)=2$, и по
примеру~\ref{ex:ample-examples}(3) $-K$ обилен. Для $\mathbb F_a$, $a\ge2$: на отрицательном сечении
$C$ ($C^2=-a$) по формуле присоединения $-K\cdot C=2+C^2=2-a\le0$ --- \emph{не} Фано.
\item \textbf{Торический критерий}: гладкое полное торическое многообразие Фано $\iff$ $\conv(u_\rho)$ ---
многогранник, а конусы веера --- ровно конусы над его гранями (теорема~\ref{thm:toric-Fano}).
\item \textbf{Накрытия.} Если $\pi\colon X\to Y$ --- двулистное накрытие, разветвлённое в гладком дивизоре
$B\in|2L|$, то $K_X=\pi^*(K_Y+L)$. Если $-(K_Y+L)$ обилен, то $X$ Фано (обратный образ обильного при
конечном морфизме обилен). Пример: $Y=\PP^2$, $B$ --- квартика, $L=\cO(2)$: $-K_X=\pi^*\cO(1)$ ---
дель Пеццо степени $2$.
\item \textbf{Раздутия.} $K_{\Bl_ZX}=\pi^*K_X+(c-1)E$ для гладкого центра $Z$ коразмерности $c$; обильность
$-K$ надо проверять заново (на кривых в $E$ и собственных прообразах).
\item \textbf{Дифференциальная геометрия}: предъявить кэлерову метрику с $\Ric>0$ (например, однородную
метрику на $G/P$, Фубини--Штуди на $\PP^n$).
\end{enumerate}

\begin{example}[Не Фано]
\begin{itemize}
\item Раздутие $\PP^2$ в трёх точках на одной прямой $\ell$: собственный прообраз $\tilde\ell$ имеет класс
$L-E_1-E_2-E_3$ и $-K\cdot\tilde\ell=3-3=0$.
\item Раздутие $\PP^2$ в 9 точках: $K^2=0$.
\item Раздутие точки $p$ и затем точки $q\in E_1$ («бесконечно близкой»): собственный прообраз $E_1$ ---
кривая класса $E_1-E_2$ с $(E_1-E_2)^2=-2$ и $-K\cdot(E_1-E_2)=0$.
\end{itemize}
\end{example}

% ============================================================
\section{Раздутия точек на проективной плоскости}\label{sec:blowup}
% ============================================================

\subsection{Раздутие точки на плоскости}

\begin{definition}
\emph{Раздутие} $\A^2$ в точке $0$:
\[
B=\Bl_0\A^2=\{((x,y),[u:v])\in\A^2\times\PP^1:\ xv=yu\},\qquad\pi\colon B\to\A^2 .
\]
\end{definition}

\begin{proposition}\label{prop:blowup-basic}
\begin{enumerate}
\item $B$ --- гладкая поверхность, покрытая двумя картами $\cong\A^2$:
$\{u\ne0\}$ с координатами $(x,t)$, $t=v/u$, где $y=xt$; и $\{v\ne0\}$ с координатами $(s,y)$, $s=u/v$,
где $x=sy$.
\item $E=\pi^{-1}(0)=\{0\}\times\PP^1$ --- \emph{исключительная кривая}; в первой карте $E=\{x=0\}$.
\item $\pi\colon B\setminus E\to\A^2\setminus0$ --- изоморфизм (обратное: $(x,y)\mapsto((x,y),[x:y])$).
\item Точки $E$ --- это прямые через $0$: собственный прообраз прямой $y=cx$ пересекает $E$ в точке $[1:c]$.
\end{enumerate}
\end{proposition}
\begin{proof}
В карте $u=1$: $v=t$, уравнение $y=xt$ задаёт график, т.\,е. $\cong\A^2$ с координатами $(x,t)$;
аналогично вторая. Остальное очевидно.
\end{proof}

\begin{figure}[ht]
\centering
\begin{tikzpicture}[scale=0.9]
% blown-up surface: lines separated
\draw[thick,red!70!black] (0,-1.6) -- (0,1.6) node[above] {$E\cong\PP^1$};
\foreach \c/\h in {-1.2/-1.2,-0.6/-0.6,0/0,0.6/0.6,1.2/1.2}{
  \draw[blue!70!black] (-2.2,\h-0.15*\c) .. controls (-1,\h) and (1,\h) .. (2.2,\h+0.15*\c);
  \fill (0,\h) circle (1.3pt);}
\node at (0,-2.3) {$\Bl_0\A^2$};
\draw[->,thick] (2.8,0) -- (4.2,0) node[midway,above] {$\pi$};
\begin{scope}[xshift=7cm]
\foreach \a in {-50,-25,0,25,50}{ \draw[blue!70!black] (\a:-2) -- (\a:2);}
\fill[red!70!black] (0,0) circle (2.3pt);
\node at (0,-2.3) {$\A^2$};
\end{scope}
\end{tikzpicture}
\caption{Раздутие: прямые через $0$ «разводятся» и пересекают $E$ в разных точках; точка $0$ заменяется на
проективную прямую направлений.}
\label{fig:blowup}
\end{figure}

\emph{Раздутие точки $p$ на гладкой поверхности $S$.} Пусть $u,v$ --- локальные параметры в $p$, регулярные
на окрестности $U$, с единственным общим нулём $p$ в $U$. Положим
\[
\Bl_pU=\overline{\{(q,[u(q):v(q)]):\ q\in U\setminus p\}}\subset U\times\PP^1
\]
и склеим с $S\setminus\{p\}$. Результат $\tilde S=\Bl_pS$ не зависит от выбора параметров и является гладкой
поверхностью с морфизмом $\pi\colon\tilde S\to S$, изоморфным вне $E=\pi^{-1}(p)\cong\PP(T_pS)$
\cite[II.4]{Sh}. Так как $\pi$ --- изоморфизм вне $E$, $\tilde S$ бирационально $S$.

\begin{example}[$\Bl_p\PP^2$]
Для $p=[1:0:0]$:
$\Bl_p\PP^2=\{([x_0:x_1:x_2],[u:v]): x_1v=x_2u\}\subset\PP^2\times\PP^1$. Вторая проекция ---
это проекция из $p$; её слои --- собственные прообразы прямых через $p$, и $\Bl_p\PP^2$ ---
$\PP^1$-расслоение над $\PP^1$, изоморфное $\mathbb F_1$. Торически: веер $\PP^2$ + луч $e_1+e_2$
(рис.~\ref{fig:star}).
\end{example}

\begin{figure}[ht]
\centering
\begin{tikzpicture}[scale=0.62]
\begin{scope}
\clip (-2.2,-2.2) rectangle (2.2,2.2);
\cfill{1}{0}{0}{1}{blue!12}\cfill{0}{1}{-1}{-1}{green!12}\cfill{-1}{-1}{1}{0}{orange!15}
\lat{2}\ray{1}{0}\ray{0}{1}\ray{-1}{-1}
\end{scope}
\node at (0,-2.8) {$\PP^2$};
\end{tikzpicture}
\quad\raisebox{1.3cm}{$\xleftarrow{\ \ \pi\ \ }$}\quad
\begin{tikzpicture}[scale=0.62]
\begin{scope}
\clip (-2.2,-2.2) rectangle (2.2,2.2);
\cfill{1}{0}{1}{1}{blue!12}\cfill{1}{1}{0}{1}{cyan!15}\cfill{0}{1}{-1}{-1}{green!12}\cfill{-1}{-1}{1}{0}{orange!15}
\lat{2}\ray{1}{0}\ray{0}{1}\ray{-1}{-1}
\draw[very thick,red!70!black,-{Stealth[length=2.2mm]}] (0,0) -- (1,1); \fill[red!70!black] (1,1) circle (1.6pt);
\end{scope}
\node[red!70!black] at (1.5,1.15) {$E$};
\node at (0,-2.8) {$\Bl_p\PP^2=\mathbb F_1$};
\end{tikzpicture}
\qquad\qquad
\begin{tikzpicture}[scale=0.55]
\fill[blue!15,draw=blue!60!black,thick] (-1,-1)--(2,-1)--(-1,2)--cycle;
\draw[red!70!black,very thick,dashed] (0,-1)--(-1,0);
\foreach \x in {-1,...,2}\foreach \y in {-1,...,2}{\fill[gray!55] (\x,\y) circle (1pt);}
\node at (0.5,-2.1) {срезание вершины};
\end{tikzpicture}
\caption{Торическое раздутие неподвижной точки --- звёздное подразбиение конуса (новый луч $e_1+e_2$
--- исключительная кривая) или, двойственно, срезание вершины многоугольника.}
\label{fig:star}
\end{figure}

\subsection{Собственные прообразы, группа Пикара, пересечения}

\begin{proposition}\label{prop:strict}
Пусть $C\subset S$ --- кривая, имеющая в $p$ кратность $m$ (т.\,е. локальное уравнение $f\in\mm_p^m\setminus
\mm_p^{m+1}$), $\tilde C$ --- её \emph{собственный прообраз} (замыкание $\pi^{-1}(C\setminus p)$). Тогда
\[
\pi^*C=\tilde C+mE,
\]
и $\tilde C\cap E$ --- это касательные направления $C$ в $p$ (нули младшей формы $f_m$).
\end{proposition}
\begin{proof}
В локальных параметрах $f=f_m(x,y)+f_{m+1}(x,y)+\dots$ ($f_k$ однородны степени $k$). В карте $y=xt$:
\[
\pi^*f=f(x,xt)=x^m\big(f_m(1,t)+xf_{m+1}(1,t)+\dots\big)=x^m\tilde f,
\]
и $\tilde f$ не делится на $x$, так как $\tilde f|_{x=0}=f_m(1,t)\not\equiv0$ (коэффициенты многочлена
$f_m(1,t)$ --- это коэффициенты формы $f_m\ne0$). Значит, $\tilde f$ --- уравнение $\tilde C$ в этой карте,
$\ord_E\pi^*f=m$, а $\tilde C\cap E=\{x=0,\ f_m(1,t)=0\}$; во второй карте (направление $x=0$) ---
аналогично с $f_m(s,1)$.
\end{proof}

\begin{example}
\begin{itemize}
\item \emph{Узловая кубика} $y^2=x^2(x+1)$ ($m=2$): $x^2t^2=x^2(x+1)$, собственный прообраз $t^2=x+1$ ---
гладкая кривая, пересекающая $E=\{x=0\}$ трансверсально в двух точках $t=\pm1$ (две касательные
$y=\pm x$). Раздутие разрешает особенность.
\item \emph{Каспидальная кубика} $y^2=x^3$: во второй карте ($x=sy$): $y^2=s^3y^3$, т.\,е. $\tilde C=\{1=s^3y\}$
--- не встречает эту карту в $E$; в первой карте: $t^2=x$ --- гладкая кривая, касающаяся $E$ в точке $t=0$
(одна касательная $y=0$ кратности 2).
\end{itemize}
\end{example}

\begin{theorem}\label{thm:blowup-Pic}
Пусть $\pi\colon\tilde S=\Bl_pS\to S$ --- раздутие точки на гладкой проективной поверхности. Тогда
\begin{enumerate}
\item $\Pic\tilde S=\pi^*\Pic S\oplus\Z E$;
\item $\pi^*D\cdot\pi^*D'=D\cdot D'$, $\pi^*D\cdot E=0$, $E^2=-1$ (т.\,е. $N_{E/\tilde S}=\cO_{\PP^1}(-1)$);
\item $K_{\tilde S}=\pi^*K_S+E$, $K_{\tilde S}^2=K_S^2-1$.
\end{enumerate}
\end{theorem}
\begin{proof}
(1) По предложению~\ref{prop:excision}: $\Z E\to\Cl\tilde S\to\Cl(\tilde S\setminus E)\to0$ точна, а
$\Cl(\tilde S\setminus E)=\Cl(S\setminus p)=\Cl S$. Отображение $\pi^*$ (обратный образ локальных
уравнений) расщепляет эту последовательность: композиция $\Cl S\xrightarrow{\pi^*}\Cl\tilde S\to
\Cl(\tilde S\setminus E)\cong\Cl S$ --- тождество. Элемент $E$ имеет бесконечный порядок, так как
$E\cdot E=-1\ne0$ (см.~(2)).

(2) Пересечения $\pi^*D\cdot\pi^*D'$ можно считать для $D,D'$, не проходящих через $p$ (двигая в классе
линейной эквивалентности: $D\sim A-B$ с очень обильными $A,B$, а гиперплоские сечения можно выбрать не
проходящими через $p$); тогда всё происходит вне $E$, где $\pi$ --- изоморфизм. По той же причине
$\pi^*D\cdot E=0$. Наконец, пусть $C$ --- гладкая кривая через $p$ (например, гиперплоское сечение). По
предложению~\ref{prop:strict} $\pi^*C=\tilde C+E$, причём $\tilde C$ пересекает $E$ трансверсально в одной точке
(в карте $C=\{y=0\}$, $\tilde C=\{t=0\}$, $E=\{x=0\}$). Поэтому
$0=\pi^*C\cdot E=\tilde C\cdot E+E^2=1+E^2$.

(3) Пусть $\omega$ --- рациональная 2-форма на $S$; около $p$ запишем $\omega=h\,dx\wedge dy$. В карте
$y=xt$: $\pi^*\omega=\pi^*h\cdot dx\wedge d(xt)=\pi^*h\cdot x\,dx\wedge dt$. Поэтому
$\ord_E\pi^*\omega=\ord_E\pi^*h+1$, а вне $E$ дивизоры совпадают; это и значит
$\divv\pi^*\omega=\pi^*\divv\omega+E$. Далее, $K_{\tilde S}^2=(\pi^*K_S)^2+2\pi^*K_S\cdot E+E^2=K_S^2-1$.
\end{proof}

\begin{corollary}\label{cor:Blr}
Пусть $X_r=\Bl_{p_1,\dots,p_r}\PP^2$ --- раздутие различных точек, $L=\pi^*H$, $E_i$ --- исключительные кривые.
Тогда
\[
\Pic X_r=\Z L\oplus\Z E_1\oplus\dots\oplus\Z E_r,\quad L^2=1,\ E_i^2=-1,\ L\cdot E_i=E_i\cdot E_j=0\ (i\ne j),
\]
\[
K_{X_r}=-3L+E_1+\dots+E_r,\qquad K_{X_r}^2=9-r .
\]
Решётка $\Pic X_r$ с формой пересечения --- это $\Z^{1,r}$ (форма $\mathrm{diag}(1,-1,\dots,-1)$).
Неприводимая кривая на $X_r$ --- либо одна из $E_i$, либо собственный прообраз $\tilde C$ плоской кривой $C$
степени $d$, имеющей кратности $m_i$ в $p_i$; его класс $dL-\sum m_iE_i$, и $\tilde C\cdot E_i=m_i\ge0$.
\end{corollary}

\begin{theorem}[{Кастельнуово, \cite[V.5.7]{Ha}; факторизация, \cite[V.5.3]{Ha}}]\label{thm:Castelnuovo}
\begin{enumerate}
\item Если $E\subset S$ --- гладкая рациональная кривая с $E^2=-1$, то существует гладкая поверхность $S'$ и
точка $p\in S'$ с $S\cong\Bl_pS'$, причём $E$ --- исключительная кривая.
\item Любой бирациональный морфизм гладких проективных поверхностей --- композиция раздутий точек.
\end{enumerate}
\end{theorem}

\subsection{Примеры}

\begin{example}[$\Bl_2\PP^2\cong\Bl_1(\PP^1\times\PP^1)$]
Проекция гладкой квадрики $Q\subset\PP^3$ из точки $q\in Q$ на плоскость не определена в $q$; после
раздутия $q$ она становится морфизмом $\Bl_qQ\to\PP^2$, стягивающим собственные прообразы двух прямых на $Q$
через $q$ (они имеют самопересечение $0-1=-1$) в две точки. Значит, $\Bl_q(\PP^1\times\PP^1)\cong
\Bl_2\PP^2$. Торически: веер $\PP^1\times\PP^1$ с добавленным лучом $-e_1-e_2$ и веер $\PP^2$ с
добавленными лучами $e_1+e_2$, $-e_1$ переводятся друг в друга элементом $\GL_2(\Z)$
(последовательности самопересечений $(0,0,-1,-1,-1)$ совпадают с точностью до циклического сдвига).
\end{example}

\begin{example}[Квадратичное преобразование Кремоны]
$\sigma\colon[x:y:z]\mapsto[yz:xz:xy]=[1/x:1/y:1/z]$ --- бирациональная инволюция $\PP^2$, не определённая в
трёх координатных точках и стягивающая три координатные прямые. На $X_3=\Bl_3\PP^2$ она становится
автоморфизмом, переставляющим шестёрку $(-1)$-кривых: $E_i\leftrightarrow L-E_j-E_k$. Торически это симметрия
$u\mapsto-u$ шестиугольного веера (рис.~\ref{fig:fans}). В $\Pic X_3$: $L\mapsto2L-E_1-E_2-E_3$.
\end{example}

\subsection{Раздутия в общем положении и поверхности дель Пеццо}

\begin{definition}
Точки $p_1,\dots,p_r\in\PP^2$ ($r\le8$) находятся в \emph{общем положении}, если никакие 3 не лежат на прямой,
никакие 6 не лежат на конике, и (при $r=8$) нет кубики через все точки, особой в одной из них.
\end{definition}

\begin{lemma}[Необходимость общего положения]
Если 3 из точек лежат на прямой $\ell$ или 6 --- на конике $Q$, то $X_r$ не Фано.
\end{lemma}
\begin{proof}
Собственный прообраз $\tilde\ell$ имеет класс $L-E_i-E_j-E_k-\dots$ и
$-K\cdot\tilde\ell=3-\#(\ell\cap\{p_i\})\le0$. Аналогично $-K\cdot\tilde Q=6-\#(Q\cap\{p_i\})\le0$.
\end{proof}

\begin{theorem}\label{thm:dP-blowup}
Если $r\le6$ и точки $p_1,\dots,p_r$ находятся в общем положении, то $X_r$ --- поверхность дель Пеццо степени
$9-r$.
\end{theorem}
\begin{proof}
По критерию Накаи--Мойшезона (теорема~\ref{thm:ample-criteria}) надо проверить $K^2=9-r>0$ (верно) и
$-K\cdot C>0$ для любой неприводимой кривой $C$. Для $C=E_i$: $-K\cdot E_i=1$.

\emph{Случай $r=6$.} Обозначим $\ell_{ij}=L-E_i-E_j$ --- класс собственного прообраза прямой $p_ip_j$
(неприводимой кривой; других точек на прямой нет). Для любого разбиения $\{1,\dots,6\}$ на три пары
\[
-K=3L-\sum_{k=1}^6E_k=\ell_{ij}+\ell_{kl}+\ell_{st}.
\]
Пусть $C$ --- собственный прообраз плоской кривой степени $d\ge1$ с кратностями $m_i$. Если $C=\ell_{ij}$, то
$-K\cdot C=3-2=1$. Иначе $C$ отлична от всех $\ell_{ij}$, и $C\cdot\ell_{ij}=d-m_i-m_j\ge0$; значит,
$-K\cdot C\ge0$, причём равенство $-K\cdot C=0$ влечёт $m_i+m_j=d$ для всех пар в каждом разбиении, т.\,е. для
всех $i\ne j$ (любая пара входит в некоторое разбиение). Отсюда $m_1=\dots=m_6=d/2$, и $d$ чётно. Рассмотрим
конику $Q$ через $p_1,\dots,p_5$ (она единственна и неприводима, так как никакие 3 точки не лежат на прямой);
её собственный прообраз имеет класс $\tilde Q=2L-E_1-\dots-E_5$ (через $p_6$ она не проходит). Тогда
\[
C\cdot\tilde Q=2d-5\cdot\tfrac d2=-\tfrac d2<0,
\]
что возможно для различных неприводимых кривых только при $C=\tilde Q$; но у $\tilde Q$ кратность в $p_6$ равна
$0\ne d/2$. Противоречие; итак, $-K\cdot C>0$.

\emph{Случай $r<6$.} Добавим точки $p_{r+1},\dots,p_6$ так, чтобы все шесть были в общем положении и новые точки
не лежали на образе $\pi(C)$ данной кривой $C\subset X_r$ (это открытые условия). Если $C=\tilde C_0$ ---
собственный прообраз плоской кривой $C_0$, то её собственный прообраз $C'$ на $X_6$ имеет те же $d,m_1,\dots,m_r$ и
$m_{r+1}=\dots=m_6=0$, поэтому $-K_{X_r}\cdot C=3d-\sum_{i\le r}m_i=-K_{X_6}\cdot C'>0$.
\end{proof}

\begin{remark}
Для $r=7,8$ доказательство аналогично, но требует перебора большего числа «плохих» классов (например, при $r=8$
--- кубика через все 8 точек, особая в одной из них: $-K\cdot C=9-2-7=0$); см. \cite[Теорема~24.4]{Ma}, \cite{De}, \cite[8.1]{Do}. При $r\ge9$
имеем $K^2\le0$, и $X_r$ не Фано.
\end{remark}

\begin{remark}[$(-1)$-кривые на $X_r$]
Класс $e$ неприводимой $(-1)$-кривой удовлетворяет $e^2=-1$, $K\cdot e=-1$ (пример~\ref{ex:adj-surfaces}).
Если $-K$ обилен, то верно и обратное: каждый такой класс содержит ровно одну неприводимую кривую (см.
\S\ref{sec:cubic} для $r=6$). Число $(-1)$-кривых на дель Пеццо $X_r$:
\[
\begin{array}{c|cccccccc}
r & 1&2&3&4&5&6&7&8\\\hline
\#(-1)\text{-кривых} & 1&3&6&10&16&27&56&240
\end{array}
\]
\end{remark}

% ============================================================
\section{Кубические поверхности и 27 прямых}\label{sec:cubic}
% ============================================================

Пусть $S=\{F=0\}\subset\PP^3$ --- гладкая кубическая поверхность.

\subsection{Кубическая поверхность --- поверхность дель Пеццо}

\begin{theorem}
$K_S=\cO_S(-1)$. Следовательно, $-K_S$ очень обилен (это гиперплоское сечение), $S$ --- поверхность дель Пеццо
степени $K_S^2=H^2\cdot S=3$, и для любой неприводимой кривой $C\subset S$: $\deg C=-K_S\cdot C$.
\end{theorem}
\begin{proof}
Формула присоединения: $K_S=(K_{\PP^3}+S)|_S=\cO(-4+3)|_S$ (пример~\ref{ex:hypersurface}).
\end{proof}

В частности, \emph{прямые на $S$ --- это в точности неприводимые кривые $C$ с $-K_S\cdot C=1$}, и для них
$C^2=-1$ (пример~\ref{ex:adj-surfaces}).

\subsection{Прямые, пересекающие данную прямую}

\begin{lemma}\label{lem:tangent-line}
Если прямая $m\subset S$ проходит через $q$, то $m\subset\overline T_qS$ (проективная касательная плоскость).
\end{lemma}
\begin{proof}
Если $q+sv\in S$ при всех $s$, то $0=\frac{d}{ds}F(q+sv)|_{s=0}=\sum v_i\partial_iF(q)$.
\end{proof}

\begin{lemma}\label{lem:double-line}
Если для плоскости $H$ сечение $H\cap S$ содержит прямую $m$ с кратностью $\ge2$, то $S$ особа в некоторой
точке $m$.
\end{lemma}
\begin{proof}
Пусть $H=\{h=0\}$, $m=\{h=\lambda=0\}$, $F|_H=\lambda^2\mu$. Тогда $F=\lambda^2\mu+hG$, $G$ --- квадратичная
форма, и $\nabla F=2\lambda\mu\nabla\lambda+\lambda^2\nabla\mu+G\nabla h+h\nabla G$. На $m$ имеем
$\nabla F=G\cdot\nabla h$, а $G|_m$ --- бинарная квадратичная форма, у которой есть нуль на $m\cong\PP^1$. В нём
$\nabla F=0$.
\end{proof}

Пусть $\ell\subset S$ --- прямая. Каждая плоскость $H\supset\ell$ высекает $H\cap S=\ell+Q_H$, где $Q_H$ ---
коника. По лемме~\ref{lem:double-line} $Q_H$ не содержит $\ell$ и не является двойной прямой, т.\,е. $Q_H$ либо
неприводима, либо есть пара различных прямых, отличных от $\ell$. Во втором случае $H$ называется
\emph{тритангенциальной плоскостью} (она касается $S$ в вершинах треугольника).

\begin{proposition}\label{prop:five-planes}
Через каждую прямую $\ell\subset S$ проходит ровно $5$ тритангенциальных плоскостей $H_1,\dots,H_5$. Поэтому
$\ell$ пересекают ровно $10$ прямых $m_i,m_i'\subset H_i$ на $S$; прямые из разных плоскостей $H_i\ne H_j$ не
пересекаются; точки $m_i\cap\ell$ ($i=1,\dots,5$) попарно различны.
\end{proposition}
\begin{proof}
Выберем координаты $(x,y,z,t)$ с $\ell=\{z=t=0\}$. Так как $F\in(z,t)$, имеем
\[
F=Ax^2+2Bxy+Cy^2+2Dx+2Ey+G,
\]
где $A,B,C$ --- линейные, $D,E$ --- квадратичные, $G$ --- кубическая формы от $(z,t)$. Плоскость через $\ell$ ---
это $H_{(z_0:t_0)}=\{(x,y,\lambda z_0,\lambda t_0)\}$ с координатами $[x:y:\lambda]$, и
\[
F|_H=\lambda\big(A x^2+2Bxy+Cy^2+2D\lambda x+2E\lambda y+G\lambda^2\big),\quad A=A(z_0,t_0),\ \dots
\]
Так что $Q_H$ --- коника с матрицей $\begin{pmatrix}A&B&D\\B&C&E\\D&E&G\end{pmatrix}$, а она вырождена при
$\Delta(z_0,t_0)=0$, где $\Delta=\det$ --- бинарная форма степени $1+1+3=5$.

Покажем: \emph{если $H=\{t=0\}$ тритангенциальна, то $t$ делит $\Delta$ ровно в первой степени.} Тогда $\Delta\not\equiv0$
(иначе все плоскости вырождены и $t^2\mid\Delta$), и все корни $\Delta$ простые: их ровно 5.

\emph{Случай 1: три прямые $\ell,L_1,L_2$ в $H$ не проходят через одну точку.} Тогда их уравнения в плоскости
$t=0$ линейно независимы, и можно выбрать координаты $x,y,z$ так, что $\ell=\{z=0\}$, $L_1=\{x=0\}$,
$L_2=\{y=0\}$: $F|_{t=0}=xyz$, т.\,е. $F=xyz+tF'$, $F'$ --- квадратичная форма. Тогда
$A=at$, $B=\frac z2+bt$, $C=ct$, $D=t\,d$, $E=t\,e$, $G=t\,g$, где $g(z,t)=F'(0,0,z,t)$. Раскрывая
определитель,
\[
\Delta=A(CG-E^2)-B(BG-ED)+D(BE-CD)\equiv-B^2G\equiv-\tfrac14F'(0,0,1,0)\,z^4\,t\pmod{t^2}.
\]
Но в точке $P=(0:0:1:0)=L_1\cap L_2\in S$: $\partial_xF=\partial_yF=\partial_zF=0$, а $\partial_tF(P)=F'(P)$.
Из гладкости $F'(P)\ne0$.

\emph{Случай 2: $\ell,L_1,L_2$ проходят через точку $P\in\ell$} (\emph{точка Эккардта}). Выберем координаты
так, что $P=(0:1:0:0)$, $L_1=\{x=t=0\}$, $L_2=\{x+z=t=0\}$; тогда $F|_{t=0}=xz(x+z)$ и $F=x^2z+xz^2+tF'$.
Теперь $A=z+at$, $B=bt$, $C=c't$, $D=\frac{z^2}2+td$, $E=te$, $G=tg$, где $c'=F'(0,1,0,0)$, и
\[
\Delta\equiv-CD^2\equiv-\tfrac{c'}4\,z^4\,t\pmod{t^2}.
\]
В точке $P$ все производные, кроме $\partial_tF(P)=c'$, равны нулю; из гладкости $c'\ne0$.

Итак, тритангенциальных плоскостей ровно пять. Всякая прямая $m\ne\ell$, пересекающая $\ell$, лежит в плоскости
$\langle\ell,m\rangle$, которая тогда тритангенциальна; это даёт 10 прямых. Пусть $m\subset H_i$, $n\subset H_j$,
$i\ne j$, и $q\in m\cap n$. Тогда $q\in H_i\cap H_j=\ell$, и прямые $\ell,m,n$ проходят через $q$ и по
лемме~\ref{lem:tangent-line} лежат в плоскости $\overline T_qS$; но они не компланарны ($m\not\subset H_j$) ---
противоречие. Последнее утверждение --- частный случай.
\end{proof}

\begin{corollary}\label{cor:skew}
На $S$, содержащей прямую, есть две непересекающиеся (скрещивающиеся) прямые: $m_1\subset H_1$ и $m_2\subset H_2$.
\end{corollary}

\begin{lemma}\label{lem:five-common}
Если $\ell_1,\ell_2\subset S$ --- скрещивающиеся прямые, то ровно 5 прямых на $S$ пересекают обе, и эти 5 прямых
попарно не пересекаются.
\end{lemma}
\begin{proof}
Прямые, пересекающие $\ell_1$, --- это $m_i,m'_i\subset H_i$, $i=1,\dots,5$ (предложение~\ref{prop:five-planes}).
Так как $\ell_2\not\subset H_i$ (иначе $\ell_2$ пересекла бы $\ell_1$), $\ell_2\cap H_i$ --- одна точка $q_i\in
S\cap H_i=\ell_1\cup m_i\cup m_i'$, причём $q_i\notin\ell_1$. Если бы $q_i\in m_i\cap m'_i$, то
$\overline T_{q_i}S\supset m_i,m_i'$, т.\,е. $\overline T_{q_i}S=H_i$, и $\ell_2\subset H_i$ по
лемме~\ref{lem:tangent-line} --- противоречие. Значит, ровно одна из $m_i,m_i'$ (обозначим её $E_i$)
пересекает $\ell_2$. Прямые $E_1,\dots,E_5$ лежат в разных $H_i$ и потому не пересекаются.
\end{proof}

\subsection{Проекция на $\ell_1\times\ell_2$: рациональность и модель раздутия}

Пусть $\ell_1,\ell_2\subset S$ скрещиваются. Выберем координаты:
$\ell_1=\{x_2=x_3=0\}$, $\ell_2=\{x_0=x_1=0\}$, и определим
\[
\varphi\colon S\to\ell_1\times\ell_2\cong\PP^1\times\PP^1,\qquad p\mapsto(a(p),b(p)),
\]
где для $p\notin\ell_1\cup\ell_2$ точки $a(p)\in\ell_1$, $b(p)\in\ell_2$ таковы, что единственная прямая через $p$,
пересекающая $\ell_1$ и $\ell_2$, --- это прямая $a(p)b(p)$: $b(p)=\langle p,\ell_1\rangle\cap\ell_2$,
$a(p)=\langle p,\ell_2\rangle\cap\ell_1$. В координатах
\[
a(p)=[x_0:x_1],\qquad b(p)=[x_2:x_3].
\]

\begin{theorem}\label{thm:projection}
\begin{enumerate}
\item $\varphi$ продолжается до морфизма $S\to\PP^1\times\PP^1$: для $p\in\ell_1$ надо положить
$a(p)=p$, $b(p)=\overline T_pS\cap\ell_2$ (и симметрично для $p\in\ell_2$).
\item $\varphi$ бирационален; обратное отображение $\psi(a,b)=$ «третья точка пересечения прямой $ab$ с $S$».
В частности, \textbf{$S$ рациональна}.
\item $\varphi$ стягивает пять прямых $E_1,\dots,E_5$ из леммы~\ref{lem:five-common} в пять различных точек
$q_i=(a_i,b_i)$ и является изоморфизмом вне них; $S\cong\Bl_{q_1,\dots,q_5}(\PP^1\times\PP^1)\cong\Bl_6\PP^2$.
\end{enumerate}
\end{theorem}
\begin{proof}
(1) Так как $F$ обращается в нуль на $\ell_1$, $F=x_2Q_2+x_3Q_3$ с квадратичными $Q_2,Q_3$. На $S$ имеем
$x_2Q_2=-x_3Q_3$, поэтому $[x_2:x_3]=[-Q_3:Q_2]$ там, где обе части определены. В точке $p\in\ell_1$:
$\nabla F(p)=(0,0,Q_2(p),Q_3(p))\ne0$ (гладкость), так что $[-Q_3:Q_2]$ определено около $p$, и
$\overline T_pS=\{Q_2(p)x_2+Q_3(p)x_3=0\}$ пересекает $\ell_2$ как раз в $[-Q_3(p):Q_2(p)]$. Значит, $b$ ---
морфизм $S\to\PP^1$; аналогично $a$.

(2) Для $a=(a_0,a_1,0,0)$, $b=(0,0,b_2,b_3)$ ограничение $F$ на прямую $\{\lambda a+\mu b\}$ делится на
$\lambda\mu$ (обе точки на $S$): $F(\lambda a+\mu b)=\lambda\mu(\alpha\lambda+\beta\mu)$, где $\alpha,\beta$ ---
биоднородные формы от $(a,b)$. Положим $\psi(a,b)=\beta a-\alpha b$ --- рациональное отображение
$\PP^1\times\PP^1\dashrightarrow S$, определённое там, где $(\alpha,\beta)\ne0$, т.\,е. прямая $ab\not\subset S$.
Точка $c=\psi(a,b)$ лежит на прямой $ab$, поэтому $\varphi(c)=(a,b)$ при $c\notin\ell_1\cup\ell_2$; значит,
$\varphi\circ\psi=\id$ на открытом множестве, и $\varphi$ бирационален. $S$ бирационально $\PP^1\times\PP^1$,
т.\,е. рациональна.

(3) Найдём слои. Пусть $\lambda=$ прямая $ab$. Если $\lambda\subset S$, то $\lambda$ пересекает $\ell_1,\ell_2$, т.\,е.
$\lambda=E_i$, и все точки $E_i$ переходят в $(a_i,b_i)$ (для $p=a_i$: $\overline T_{a_i}S=\langle\ell_1,E_i\rangle$
пересекает $\ell_2$ в $b_i$). Если $\lambda\not\subset S$, то $\varphi^{-1}(a,b)=\{c\}$: точки $p\notin\ell_1\cup\ell_2$
со $\varphi(p)=(a,b)$ лежат на $\lambda$, т.\,е. $p=c$; а точка $p=a\in\ell_1$ попадает в слой, только если
$b\in\overline T_aS$, т.\,е. $\lambda$ касается $S$ в $a$, т.\,е. $c=a$ (двойной корень); аналогично для $\ell_2$.

Итак, $\varphi$ биективен над дополнением к точкам $q_i=\varphi(E_i)$, которые попарно различны: $a_i=E_i\cap\ell_1$
различны по предложению~\ref{prop:five-planes}. Биективный бирациональный морфизм на гладкое (нормальное)
многообразие --- изоморфизм (главная теорема Зарисского \cite[III.9]{Sh}, \cite[III.11.4]{Ha}). По
теореме~\ref{thm:Castelnuovo}(2) $\varphi$ --- композиция $k$ раздутий точек; исключительное множество композиции
имеет ровно $k$ неприводимых компонент, и если бы некоторое раздутие было в точке, лежащей на исключительной
кривой предыдущих, то над одной точкой лежали бы $\ge2$ пересекающиеся компоненты. Поскольку $E_1,\dots,E_5$ не
пересекаются и лежат над разными точками, $\varphi$ --- раздутие пяти различных точек $q_1,\dots,q_5$.
Наконец, $q_j$ ($j\ge2$) не лежат на прямых $\{a_1\}\times\PP^1$, $\PP^1\times\{b_1\}$ (все $a_i$ различны, все $b_i$
различны), и по примеру из \S\ref{sec:blowup} $\Bl_{q_1}(\PP^1\times\PP^1)\cong\Bl_2\PP^2$, причём остальные 4
точки переходят в 4 точки $\PP^2$ вне двух центров. Итого $S\cong\Bl_6\PP^2$.
\end{proof}

\subsection{27 прямых}

\begin{theorem}\label{thm:27}
Гладкая кубическая поверхность $S$, содержащая прямую (а в п.~\ref{ss:existence} мы увидим, что это любая гладкая
кубика), содержит ровно 27 прямых. В модели $S\cong X_6=\Bl_{p_1,\dots,p_6}\PP^2$ точки $p_i$ находятся в общем
положении, вложение $S\subset\PP^3$ задаётся линейной системой $|-K|=|3L-\sum E_i|$ кубик через $p_1,\dots,p_6$, и
прямые --- это
\[
\underbrace{E_i}_{6\ \text{штук}},\qquad
\underbrace{\ell_{ij}=L-E_i-E_j}_{15:\ \text{прямые } p_ip_j},\qquad
\underbrace{Q_j=2L-\textstyle\sum_{k\ne j}E_k}_{6:\ \text{коники через 5 точек}}.
\]
\end{theorem}
\begin{proof}
По теореме~\ref{thm:projection} $S\cong X_6$, и $\cO_S(1)=-K_S=3L-\sum E_i$ (корректно: $K_S$ определён
внутренне, а $K_{X_6}$ вычислен в следствии~\ref{cor:Blr}). Более того, вложение $S\subset\PP^3$ задаётся
\emph{полной} линейной системой $|-K_S|$: из точной последовательности $0\to\cO_{\PP^3}(-2)\to\cO_{\PP^3}(1)\to
\cO_S(1)\to0$ и $H^1(\PP^3,\cO(-2))=0$ получаем $h^0(\cO_S(1))=4$, а сечения $3L-\sum E_i$ на $X_6$ --- это
кубики через $p_1,\dots,p_6$.

\emph{Общее положение.} Так как $-K_S$ обилен, $-K\cdot C>0$ для всех кривых, и по лемме о необходимости общего
положения (\S\ref{sec:blowup}) никакие 3 точки не лежат на прямой и никакие 6 --- на конике.

\emph{Классы.} Прямая $C\subset S$ --- неприводимая кривая с $-K\cdot C=1$ и $C^2=-1$. Запишем
$C=dL-\sum m_iE_i$. Условия:
\[
3d-\sum_{i=1}^6m_i=1,\qquad d^2-\sum_{i=1}^6m_i^2=-1.
\]
По неравенству Коши--Буняковского $(3d-1)^2=(\sum m_i)^2\le6\sum m_i^2=6(d^2+1)$, т.\,е. $3d^2-6d-5\le0$,
откуда $d\in\{0,1,2\}$.
\begin{itemize}
\item $d=0$: $\sum m_i=-1$, $\sum m_i^2=1$ $\Rightarrow$ одно $m_i=-1$: класс $E_i$.
\item $d=1$: $\sum m_i=2$, $\sum m_i^2=2$ $\Rightarrow$ два $m_i=1$: класс $L-E_i-E_j$.
\item $d=2$: $\sum m_i=5$, $\sum m_i^2=5$ $\Rightarrow$ пять $m_i=1$: класс $2L-\sum_{k\ne j}E_k$.
\end{itemize}
Всего $6+15+6=27$ классов.

\emph{В каждом классе ровно одна прямая.} Не более одной: две различные неприводимые кривые $C\ne C'$ одного класса
$e$ дали бы $e^2=C\cdot C'\ge0$, а $e^2=-1$. Не менее одной: $E_i$ --- исключительная кривая; $\ell_{ij}$ ---
собственный прообраз прямой $p_ip_j$ (на ней нет других точек $p_k$, так что класс ровно $L-E_i-E_j$); $Q_j$ ---
собственный прообраз коники через 5 точек $p_k$, $k\ne j$ (она единственна, неприводима, так как никакие 3 точки
не на прямой, и не проходит через $p_j$). Все они неприводимы и имеют $-K\cdot C=1$, т.\,е. являются прямыми в
$\PP^3$.
\end{proof}

\begin{figure}[ht]
\centering
\begin{tikzpicture}[scale=1.0]
\draw[thick,green!50!black] (0,0) ellipse (2.4 and 1.4);
\foreach \a/\n in {15/1,90/2,155/3,215/4,300/5}{
  \fill ({2.4*cos(\a)},{1.4*sin(\a)}) circle (2.2pt);
  \node at ({2.75*cos(\a)},{1.7*sin(\a)}) {$p_{\n}$};}
\coordinate (p6) at (0.5,0.6);
\fill (p6) circle (2.2pt); \node[below] at (p6) {$p_6$};
\coordinate (p1) at ({2.4*cos(15)},{1.4*sin(15)});
\draw[thick,blue!70!black,shorten <=-2.2cm,shorten >=-1.2cm] (p6) -- (p1);
\node[blue!70!black] at (3.7,0.95) {$\ell_{16}=L-E_1-E_6$};
\node[green!50!black] at (-1.3,-1.85) {$Q_6=2L-E_1-\dots-E_5$};
\end{tikzpicture}
\caption{Шесть точек общего положения на $\PP^2$. Прямые на кубической поверхности: 6 исключительных кривых $E_i$,
15 собственных прообразов прямых $p_ip_j$, 6 собственных прообразов коник через 5 точек.}
\label{fig:sixpoints}
\end{figure}

\begin{proposition}[Конфигурация 27 прямых]
Индексы пересечения:
\[
E_i\cdot E_j=0,\quad E_i\cdot\ell_{jk}=\begin{cases}1,&i\in\{j,k\}\\0\end{cases},\quad
E_i\cdot Q_j=\begin{cases}1,&i\ne j\\0,&i=j\end{cases},
\]
\[
\ell_{ij}\cdot\ell_{kl}=\begin{cases}1,&\{i,j\}\cap\{k,l\}=\varnothing\\0,&|\{i,j\}\cap\{k,l\}|=1\end{cases},\quad
\ell_{ij}\cdot Q_k=\begin{cases}1,&k\in\{i,j\}\\0\end{cases},\quad Q_i\cdot Q_j=0\ (i\ne j).
\]
Каждая прямая пересекает ровно 10 других (ср. предложение~\ref{prop:five-planes}): например, $E_1$ пересекает
$\ell_{1j}$ и $Q_j$, $j\ne1$. Наборы $\{E_1,\dots,E_6\}$ и $\{Q_1,\dots,Q_6\}$ образуют \emph{двойную шестёрку
Шлефли}: прямые внутри каждой шестёрки попарно скрещиваются, а $E_i$ пересекает $Q_j$ ровно при $i\ne j$.
Тритангенциальных плоскостей $27\cdot5/3=45$: это тройки $\{E_i,\ell_{ij},Q_j\}$ ($30$ штук) и
$\{\ell_{ij},\ell_{kl},\ell_{st}\}$ для разбиений на пары ($15$ штук).
\end{proposition}
\begin{proof}
Прямое вычисление в $\Z^{1,6}$; например, $Q_i\cdot Q_j=4-\#(\{1..6\}\setminus\{i,j\})=0$.
\end{proof}

\begin{remark}[Симметрия $E_6$]
Классы $\alpha\in\Pic S$ с $\alpha^2=-2$, $K\cdot\alpha=0$ образуют систему корней типа $E_6$ в $K^\perp$:
$E_i-E_j$ ($30$), $\pm(L-E_i-E_j-E_k)$ ($40$), $\pm(2L-\sum E_i)$ ($2$) --- всего 72 корня. Группа Вейля
$W(E_6)$ порядка $51840$ действует на 27 классах, сохраняя индексы пересечения, и является группой симметрий
конфигурации 27 прямых. Разным выборам шестёрки попарно скрещивающихся прямых соответствуют разные
стягивания $S\to\PP^2$ (всего их $72$).
\end{remark}

\begin{example}[Кубика Ферма]
$S=\{x^3+y^3+z^3+w^3=0\}$ гладкая ($\nabla F=3(x^2,y^2,z^2,w^2)$). Она содержит прямые
\[
x=-\varepsilon y,\ z=-\varepsilon'w\qquad(\varepsilon^3=\varepsilon'^3=1),
\]
и аналогичные для разбиений $\{x,z\}\{y,w\}$ и $\{x,w\}\{y,z\}$: $3\cdot9=27$ различных прямых. По
теореме~\ref{thm:27} других прямых на ней нет.
\end{example}

\subsection{Каждая гладкая кубическая поверхность содержит прямую}\label{ss:existence}

\begin{theorem}
Каждая кубическая поверхность в $\PP^3$ содержит прямую. Следовательно, каждая гладкая кубическая поверхность
содержит ровно 27 прямых, изоморфна $\Bl_6\PP^2$ (6 точек в общем положении) и рациональна.
\end{theorem}
\begin{proof}
Кубические поверхности образуют $\PP^{19}=\PP(\Sym^3(\C^4)^*)$. Рассмотрим
$\Phi=\{(\ell,S):\ell\subset S\}\subset\Gr(2,4)\times\PP^{19}$ --- замкнутое подмножество. Для фиксированной прямой
$\ell$ ограничение $H^0(\PP^3,\cO(3))\to H^0(\ell,\cO(3))\cong\C^4$ сюръективно, поэтому кубики, содержащие $\ell$,
образуют $\PP^{15}$. Значит, $\Phi$ --- расслоение над $\Gr(2,4)$ (размерности 4) со слоем $\PP^{15}$; оно
неприводимо и $\dim\Phi=19$ \cite[I.6.3, Теорема~8]{Sh}. Образ $p(\Phi)$ проекции $p\colon\Phi\to\PP^{19}$ замкнут
(теорема~\ref{thm:complete}) и неприводим. Слой над кубикой Ферма конечен (27 точек). По теореме о размерности слоёв
каждая компонента непустого слоя имеет размерность $\ge\dim\Phi-\dim p(\Phi)$; значит, $\dim p(\Phi)=19$ и
$p(\Phi)=\PP^{19}$. Остальное --- теоремы~\ref{thm:projection} и~\ref{thm:27}.
\end{proof}

\subsection{Обратная конструкция}

\begin{theorem}[{\cite[V.4.6]{Ha}, \cite[IV]{Be}}]\label{thm:converse}
Пусть $p_1,\dots,p_6\in\PP^2$ в общем положении. Пространство $V$ кубик через $p_1,\dots,p_6$ четырёхмерно, и
отображение $X_6\to\PP(V^*)=\PP^3$, заданное линейной системой $|3L-\sum E_i|=|-K_{X_6}|$, --- вложение, образ
которого --- гладкая кубическая поверхность.
\end{theorem}
\begin{proof}[Набросок]
\emph{$\dim V=4$:} шесть точек накладывают независимые условия на кубики ($\dim=10$), так как для каждого $i$ есть
кубика через $p_k$, $k\ne i$, не проходящая через $p_i$ --- произведение коники через пять точек $p_k$, $k\ne i$
(она не проходит через $p_i$) и прямой, не проходящей через $p_i$. \emph{Вложение:} нужно проверить, что сечения
$V$ разделяют точки и касательные векторы на $X_6$; это делается с помощью приводимых кубик вида «прямая + коника»
и «три прямые» через данные точки, разбором случаев (точки на $E_i$, на $\ell_{ij}$ и т.\,д.), и использует общее
положение. \emph{Степень образа} $(-K)^2=9-6=3$; гладкость --- так как это вложение.
\end{proof}

\begin{remark}[Модули]
Шестёрки точек на $\PP^2$ по модулю $\PGL_3$ зависят от $12-8=4$ параметров, а кубические поверхности по модулю
$\PGL_4$ --- от $19-15=4$ параметров: обе конструкции дают одно и то же четырёхмерное пространство модулей.
\end{remark}

\subsection{Рациональность: итоги и комментарии}

\begin{enumerate}
\item Гладкая кубическая поверхность над $\C$ рациональна, причём двумя способами:
\emph{явно} --- через бирациональное отображение $S\dashrightarrow\ell_1\times\ell_2$ на произведение двух
скрещивающихся прямых (теорема~\ref{thm:projection}), обратное к которому --- «третья точка пересечения прямой
$ab$ с $S$»; и \emph{структурно} --- $S\cong\Bl_6\PP^2$. При этом рациональное отображение $\PP^2\dashrightarrow
S\subset\PP^3$ задаётся кубиками через 6 точек.
\item Особая кубика с изолированной двойной точкой $p$ также рациональна: проекция из $p$ бирационально отображает
$S$ на $\PP^2$ (общая прямая через $p$ пересекает $S$ в $p$ с кратностью 2 и ещё в одной точке).
\item Над незамкнутыми полями ситуация тоньше: кубика с рациональной точкой унирациональна, но не обязательно
рациональна (Сегре, Манин \cite{Ma}); критерий --- в терминах действия группы Галуа на 27 прямых.
\item В старших размерностях: гладкие кубические трёхмерные многообразия нерациональны (Клеменс--Гриффитс, через
промежуточный якобиан), а вопрос о рациональности общей кубики в $\PP^5$ долго оставался открытым (в 2025~г. Кацарков, Концевич,
Пантев и Ю анонсировали доказательство нерациональности очень общей кубики; работа проходит проверку).
\end{enumerate}


\begin{thebibliography}{99}
\bibitem[Be]{Be} A.~Beauville, \emph{Complex algebraic surfaces}, LMS Student Texts 34, Cambridge, 1996.
\bibitem[CLS]{CLS} D.~Cox, J.~Little, H.~Schenck, \emph{Toric varieties}, GSM 124, AMS, 2011.
\bibitem[Da]{Da} В.\,И.~Данилов, \emph{Геометрия торических многообразий}, УМН 33:2 (1978), 85--134.
\bibitem[De]{De} M.~Demazure, \emph{Surfaces de del Pezzo I--V}, in: S\'eminaire sur les singularit\'es des surfaces, LNM 777, Springer, 1980.
\bibitem[Do]{Do} I.~Dolgachev, \emph{Classical algebraic geometry: a modern view}, Cambridge, 2012.
\bibitem[Fu]{Fu} W.~Fulton, \emph{Introduction to toric varieties}, Princeton, 1993.
\bibitem[GH]{GH} Ф.~Гриффитс, Дж.~Харрис, \emph{Принципы алгебраической геометрии}, т.~1--2, М.: Мир, 1982.
\bibitem[Ha]{Ha} Р.~Хартсхорн, \emph{Алгебраическая геометрия}, М.: Мир, 1981.
\bibitem[Hu]{Hu} D.~Huybrechts, \emph{Complex geometry. An introduction}, Springer, 2005.
\bibitem[IP]{IP} V.\,A.~Iskovskikh, Yu.\,G.~Prokhorov, \emph{Fano varieties}, Encyclopaedia Math. Sci. 47, Springer, 1999.
\bibitem[KSC]{KSC} J.~Koll\'ar, K.~Smith, A.~Corti, \emph{Rational and nearly rational varieties}, Cambridge, 2004.
\bibitem[La]{La} R.~Lazarsfeld, \emph{Positivity in algebraic geometry I}, Springer, 2004.
\bibitem[Ma]{Ma} Ю.\,И.~Манин, \emph{Кубические формы}, М.: Наука, 1972.
\bibitem[Mu]{Mu} Д.~Мамфорд, \emph{Алгебраическая геометрия I. Комплексные проективные многообразия}, М.: Мир, 1979.
\bibitem[Re]{Re} M.~Reid, \emph{Undergraduate algebraic geometry}, LMS Student Texts 12, Cambridge, 1988.
\bibitem[Se]{Se} J.-P.~Serre, \emph{G\'eom\'etrie alg\'ebrique et g\'eom\'etrie analytique}, Ann. Inst. Fourier 6 (1956), 1--42.
\bibitem[Sh]{Sh} И.\,Р.~Шафаревич, \emph{Основы алгебраической геометрии}, т.~1--2, 3-е изд., М.: МЦНМО, 2007.
\end{thebibliography}
\end{document}
